% =================================================================
%  pccx v002 Instruction Set Architecture — Software Developer's Manual
%
%  Authoring entry point for the offline, signoff-ready PDF mirror of
%  the v002 ISA.  Styled after the Intel 64 and IA-32 Architectures
%  Software Developer's Manual (per-instruction reference blocks,
%  bit-level layouts, Operation pseudocode, Flags / Exceptions) and
%  the NVIDIA CUDA C Programming Guide (heavy TikZ architectural
%  diagrams, memory hierarchy, execution model).
%
%  Compile:  xelatex main.tex   (repeat 2x for TOC resolution)
%  Build:    bash tools/build_isa_pdf.sh
% =================================================================
\documentclass[11pt,a4paper,twoside,openany]{report}

% ---- Korean support (XeLaTeX) — enabled if the system has a Korean
% font available.  Kept optional so CI without the font still builds.
% \usepackage[xetex]{kotex}

% ---------- Typography ----------
\usepackage{fontspec}
\usepackage{microtype}

% Prefer Courier Prime when present under FONT/; otherwise fall back
% to Latin Modern.  CI passes without the house font.
\IfFileExists{FONT/CourierPrime-Regular.ttf}{%
  \setmainfont{CourierPrime-Regular.ttf}[
    Path           = FONT/,
    BoldFont       = CourierPrime-Bold.ttf,
    ItalicFont     = CourierPrime-Italic.ttf,
    BoldItalicFont = CourierPrime-BoldItalic.ttf
  ]%
}{%
  \setmainfont{Latin Modern Roman}%
  \setmonofont{Latin Modern Mono}%
  \setsansfont{Latin Modern Sans}%
}

% ---------- Core packages ----------
\usepackage{colortbl}
\usepackage{booktabs}
\usepackage{array}
\usepackage{makecell}
\usepackage{tabularx}
\usepackage{longtable}
\usepackage{multirow}
\usepackage{graphicx}
\usepackage{amsmath}
\usepackage{amssymb}
\usepackage{amsfonts}
\usepackage{etoolbox}
\usepackage{listings}
\usepackage{xcolor}
\usepackage{geometry}
\usepackage{fancyhdr}
\usepackage[explicit]{titlesec}
\usepackage{enumitem}
\usepackage{tocloft}
\usepackage{caption}
\usepackage{subcaption}

% ---------- Visualization packages ----------
\usepackage{tikz}
\usetikzlibrary{
  arrows.meta,
  calc,
  positioning,
  shapes.geometric,
  shapes.multipart,
  shapes.symbols,
  shapes.misc,
  fit,
  matrix,
  backgrounds,
  decorations.pathreplacing,
  decorations.markings,
  decorations.pathmorphing,
  patterns,
  patterns.meta,
  chains,
  scopes
}
\usepackage{pgfplots}
\pgfplotsset{compat=1.18}

% ---------- Bit-level encoding diagrams (Intel-style) ----------
\usepackage{bytefield}

% ---------- Boxed instruction reference blocks (Intel SDM style) ----
\usepackage[framemethod=TikZ]{mdframed}
\usepackage{tcolorbox}
\tcbuselibrary{skins,breakable}

% ---------- Tables larger in inner text density ----------
\AtBeginEnvironment{tabular}{\fontsize{8.5pt}{11pt}\selectfont}
\AtBeginEnvironment{longtable}{\fontsize{8.5pt}{11pt}\selectfont}
\AtBeginEnvironment{tabularx}{\fontsize{8.5pt}{11pt}\selectfont}

% ---------- Page geometry ----------
\geometry{top=2.3cm, bottom=2.3cm, left=2cm, right=2cm,
          headheight=14pt, headsep=20pt, footskip=30pt}

% Allow a small budget of extra stretch to break lines containing long
% \texttt paths and URLs that hyphenation alone cannot handle.  Keeps
% CI quiet without visibly ugly justification on normal prose.
\setlength{\emergencystretch}{3em}

% ---------- Colour palette ----------
\definecolor{oceanblue}{HTML}{006994}
\definecolor{OCEANBLUE}{HTML}{006994}
\definecolor{deepnavy}{HTML}{003E5C}
\definecolor{skytint}{HTML}{E8F4F9}
\definecolor{graytext}{HTML}{555555}
\definecolor{pureblack}{HTML}{000000}
\definecolor{codebg}{HTML}{F6F8FA}
\definecolor{accentred}{HTML}{B22222}
\definecolor{accentgreen}{HTML}{2E7D32}
\definecolor{accentorange}{HTML}{CC6600}
\definecolor{accentpurple}{HTML}{6A4C93}
\definecolor{fadegray}{HTML}{B0B0B0}

% ---------- Hyperlinks ----------
\usepackage{hyperref}
\hypersetup{
    colorlinks = true,
    linkcolor  = oceanblue,
    urlcolor   = oceanblue,
    citecolor  = oceanblue,
    pdftitle   = {pccx v002 Instruction Set Architecture — Software Developer's Manual},
    pdfauthor  = {pccx Project},
    pdfsubject = {NPU Instruction Set Architecture Reference Manual},
    pdfkeywords= {pccx, NPU, ISA, KV260, DSP48E2, systolic array, GEMM, GEMV, LLM, Transformer, CUDA, Intel SDM}
}

% ---------- TOC styling ----------
\renewcommand{\contentsname}{Table of Contents}
\renewcommand{\cftchapfont}{\sffamily\bfseries\color{oceanblue}}
\renewcommand{\cftchappagefont}{\sffamily\bfseries\color{oceanblue}}
\setlength{\cftbeforechapskip}{6pt}

% ---------- Chapter / Section layout ----------
\titleformat{\chapter}[block]
  {\normalfont\large\sffamily\bfseries\color{oceanblue}\raggedleft}
  {}
  {0pt}
  {\MakeUppercase{Chapter \thechapter}\\[-0.4ex]\huge #1}
  [\vspace{0.5ex}\color{oceanblue}\titlerule]

\titleformat{name=\chapter,numberless}[block]
  {\normalfont\large\sffamily\bfseries\color{oceanblue}\raggedleft}
  {}
  {0pt}
  {\huge #1}
  [\vspace{0.5ex}\color{oceanblue}\titlerule]

\titlespacing*{\chapter}{0pt}{-45pt}{30pt}
\titlespacing*{name=\chapter,numberless}{0pt}{-45pt}{30pt}

\titleformat{\section}
  {\normalfont\Large\sffamily\bfseries\color{oceanblue}}
  {\thesection}{1em}{#1}

\titleformat{\subsection}
  {\normalfont\large\sffamily\bfseries\color{oceanblue}}
  {\thesubsection}{1em}{#1}

\titleformat{\subsubsection}
  {\normalfont\normalsize\sffamily\bfseries\color{deepnavy}}
  {\thesubsubsection}{1em}{#1}

% ---------- Header / Footer ----------
\pagestyle{fancy}
\fancyhf{}
\renewcommand{\headrulewidth}{0.4pt}
\renewcommand{\headrule}{\color{oceanblue}\hrule height 0.4pt}
\renewcommand{\chaptermark}[1]{\markboth{#1}{}}
\renewcommand{\sectionmark}[1]{\markright{\thesection.\ #1}}

\fancyhead[LE]{\sffamily\color{oceanblue}\footnotesize\leftmark}
\fancyhead[RE]{\sffamily\color{graytext}\scriptsize pccx v002 ISA SDM}
\fancyhead[LO]{\sffamily\color{graytext}\scriptsize pccx v002 ISA SDM}
\fancyhead[RO]{\sffamily\color{oceanblue}\footnotesize\rightmark}
\fancyfoot[LE]{\sffamily\mdseries\color{pureblack}\scriptsize\thepage\hspace{1em}\color{graytext}\rightmark}
\fancyfoot[RO]{\sffamily\mdseries\color{graytext}\scriptsize\rightmark\hspace{1em}\color{pureblack}\thepage}

\fancypagestyle{plain}{
  \fancyhf{}
  \renewcommand{\headrulewidth}{0pt}
  \fancyfoot[LE]{\sffamily\color{pureblack}\small\thepage}
  \fancyfoot[RO]{\sffamily\color{pureblack}\small\thepage}
}

\setlength{\parindent}{0pt}
\setlength{\parskip}{0.8em}

% ---------- Code listings ----------
\lstdefinestyle{sv}{
  basicstyle       = \ttfamily\footnotesize,
  keywordstyle     = \bfseries\color{oceanblue},
  commentstyle     = \itshape\color{graytext},
  stringstyle      = \color{accentgreen},
  backgroundcolor  = \color{codebg},
  frame            = single,
  framesep         = 3pt,
  rulecolor        = \color{graytext!40},
  breaklines       = true,
  showstringspaces = false,
  columns          = fullflexible,
  keepspaces       = true,
  morekeywords={typedef,enum,struct,packed,module,endmodule,input,output,logic,
                always_ff,always_comb,posedge,negedge,unique,case,endcase,if,else,
                begin,end,assign,wire,reg,for,while,parameter,localparam}
}
\lstdefinestyle{pseudo}{
  basicstyle       = \ttfamily\footnotesize,
  keywordstyle     = \bfseries\color{oceanblue},
  commentstyle     = \itshape\color{graytext},
  backgroundcolor  = \color{skytint!30},
  frame            = leftline,
  framesep         = 4pt,
  rulecolor        = \color{oceanblue!60},
  breaklines       = true,
  showstringspaces = false,
  columns          = fullflexible,
  keepspaces       = true,
  morekeywords={IF,THEN,ELSE,ELSIF,ENDIF,WHILE,DO,ENDWHILE,FOR,TO,ENDFOR,
                BEGIN,END,RETURN,READ,WRITE,SIGNAL,WAIT,DISPATCH,
                DEC,ENC,CASE,OF,OTHERWISE}
}
\lstset{style=sv}

% ---------- Intel-SDM-style Instruction Reference environment ----
% \begin{InstructionRef}{MNEMONIC}{One-line description}
%   ... body ...
% \end{InstructionRef}
\newtcolorbox{InstructionRef}[2]{
  enhanced,
  breakable,
  colback   = white,
  colframe  = oceanblue,
  colbacktitle = oceanblue,
  coltitle  = white,
  fonttitle = \sffamily\bfseries\Large,
  title     = {\MakeUppercase{#1}\hfill\mdseries\normalsize #2},
  sharp corners = downhill,
  boxrule   = 1.2pt,
  leftrule  = 1.2pt,
  rightrule = 1.2pt,
  toprule   = 1.2pt,
  bottomrule= 1.2pt,
  top       = 6pt,
  bottom    = 6pt,
  left      = 8pt,
  right     = 8pt,
  before skip = 10pt,
  after skip  = 10pt,
}

% ---------- Intel-style "sub-block" inside InstructionRef ----------
\newcommand{\RefSection}[1]{%
  \par\medskip
  {\sffamily\bfseries\color{oceanblue}\normalsize\MakeUppercase{#1}}
  \par\vspace{-0.4em}
  {\color{oceanblue}\rule{\linewidth}{0.4pt}}\par\smallskip
}

% ---------- Compact "side-note" callout ----------
\newtcolorbox{SideNote}[1]{
  enhanced,
  colback      = skytint,
  colframe     = deepnavy,
  colbacktitle = deepnavy,
  coltitle     = white,
  sharp corners,
  boxrule      = 0.6pt,
  title        = {#1},
  fonttitle    = \sffamily\bfseries\small,
  fontupper    = \small,
  top=4pt, bottom=4pt, left=6pt, right=6pt,
  before skip=6pt, after skip=6pt,
}

% ---------- "Warning" / "Caution" (compliance) ----------
\newtcolorbox{CautionBox}{
  enhanced,
  colback   = accentred!5,
  colframe  = accentred,
  sharp corners,
  boxrule   = 0.8pt,
  title     = {Caution},
  coltitle  = white,
  colbacktitle = accentred,
  fonttitle = \sffamily\bfseries\small,
  fontupper = \small,
  top=4pt, bottom=4pt, left=6pt, right=6pt,
  before skip=6pt, after skip=6pt,
}

% ---------- Bit-field width — bytefield takes the width per environment
% via the [bitwidth=...] option, so nothing is set globally here.
% ----------

% ---------- Helper: render \uop with a math mu so the symbol does not
% drop out when the active main font lacks Greek lowercase coverage
% (e.g. Latin Modern Roman, our CI fallback).  xspace lets us write
% "\uop is" without losing the trailing space. ----------
\usepackage{xspace}
\providecommand{\uop}{\ensuremath{\mu}\text{op}\xspace}
\providecommand{\uops}{\ensuremath{\mu}\text{ops}\xspace}


% ---------- TikZ common styles ----------
% Intel-SDM / NVIDIA-CUDA minimalism: thin black borders, light fills, black arrows.
\tikzset{
  arr/.style          = {-Stealth, thin, draw=black},
  arrfat/.style       = {-{Stealth[length=2.5mm,width=2mm]}, line width=0.8pt, draw=black},
  arrdim/.style       = {-Stealth, thin, draw=black!55},
  bus/.style          = {line width=1.5pt, draw=black!75},
  ctrl/.style         = {-Stealth, thin, draw=black, dashed},
  trunk/.style        = {thin, draw=black, fill=none},
  databg/.style       = {fill=black!4},
  wirelabel/.style    = {font=\footnotesize\sffamily, fill=white, inner sep=1.5pt},
  blocknode/.style    = {rectangle, draw=black, thin,
                          fill=oceanblue!8, minimum width=2.4cm, minimum height=0.8cm,
                          font=\sffamily\small, align=center},
  hotblock/.style     = {rectangle, draw=black, thin,
                          fill=accentred!8, minimum width=2.4cm, minimum height=0.8cm,
                          font=\sffamily\small, align=center},
  ram/.style          = {rectangle, draw=black, thin,
                          fill=accentpurple!8, minimum width=2.0cm, minimum height=0.7cm,
                          font=\sffamily\small, align=center},
  stage/.style        = {rectangle, draw=black, thin,
                          fill=oceanblue!8, minimum width=1.6cm, minimum height=1.0cm,
                          font=\sffamily\footnotesize, align=center},
  scheduler/.style    = {rectangle, draw=black!55, thin, dashed,
                          fill=white, minimum width=3cm, minimum height=0.9cm,
                          font=\sffamily\small, align=center},
  port/.style         = {circle, draw=black, thin, fill=black!10,
                          minimum size=0.35cm, inner sep=0pt,
                          font=\sffamily\scriptsize},
  domainlabel/.style  = {font=\sffamily\scriptsize, align=center},
  everylabel/.style   = {font=\sffamily\footnotesize, align=center},
  tightarr/.style     = {-{Stealth[length=1.5mm,width=1.2mm]}, thin, draw=black},
}

% =================================================================
\begin{document}
\color{pureblack}

% =================================================================
%  Title page
% =================================================================
\begin{titlepage}
\thispagestyle{empty}
\begin{tikzpicture}[remember picture, overlay]
  % Top accent band
  \fill[oceanblue] (current page.north west) rectangle ++(\paperwidth, -0.9cm);
  % Side rule
  \draw[oceanblue, line width=3pt] ([xshift=2cm, yshift=-2.4cm]current page.north west) -- ++(0, -22cm);
\end{tikzpicture}

\vspace*{3.5cm}
\begin{flushleft}
\hspace{2cm}%
\begin{minipage}{0.82\textwidth}
{\sffamily\fontsize{44pt}{48pt}\selectfont\bfseries\color{oceanblue} pccx\par}
\vspace{0.2cm}
{\sffamily\fontsize{22pt}{26pt}\selectfont\color{deepnavy} Parallel Compute Core eXecutor\par}
\vspace{2cm}
{\sffamily\fontsize{28pt}{32pt}\selectfont\bfseries\color{pureblack} Instruction Set Architecture\par}
\vspace{0.3cm}
{\sffamily\fontsize{22pt}{26pt}\selectfont\color{pureblack} Software Developer's Manual\par}
\vspace{0.5cm}
{\sffamily\Large\color{graytext} Volume 1 --- Basic Architecture\par}
{\sffamily\Large\color{graytext} Volume 2 --- Instruction Set Reference\par}
{\sffamily\Large\color{graytext} Volume 3 --- System Programming Guide\par}

\vspace{4cm}
{\sffamily\normalsize\color{graytext}
Edge-FPGA NPU for Transformer-based LLM decoding\\[0.2em]
Target platform: Xilinx Kria KV260 (Zynq UltraScale+ ZU5EV)\\[0.2em]
Revision: v002 \quad\textbullet\quad Order Number: pccx-v002-SDM\par
}
\end{minipage}
\end{flushleft}
\vfill
\begin{tikzpicture}[remember picture, overlay]
  \fill[oceanblue] (current page.south west) rectangle ++(\paperwidth, 0.9cm);
  \node[anchor=south west, text=white, font=\sffamily\small]
    at ([xshift=2cm, yshift=0.3cm]current page.south west)
    {Document Repository: \texttt{https://github.com/hwkim-dev/pccx}};
  \node[anchor=south east, text=white, font=\sffamily\small]
    at ([xshift=-2cm, yshift=0.3cm]current page.south east)
    {Apache License 2.0};
\end{tikzpicture}
\end{titlepage}

% =================================================================
%  Front matter
% =================================================================
\pagenumbering{roman}

% ---------- Preface ----------
\chapter*{Preface}
\addcontentsline{toc}{chapter}{Preface}
\markboth{Preface}{}

This manual is the authoritative reference for the
\textbf{\color{oceanblue}pccx v002 Instruction Set Architecture}, a
fixed-length 64-bit VLIW-style ISA that drives a dedicated Neural
Processing Unit (NPU) built to accelerate autoregressive decoding of
Transformer-based Large Language Models on the Xilinx Kria KV260
platform (Zynq UltraScale+ ZU5EV).

The manual is organised as three conceptual volumes bound into a
single book: Volume 1 describes the \textbf{basic architecture} of
the NPU and its interaction with the host processing system; Volume 2
is the per-instruction reference written in the classical Intel SDM
format (Opcode table, Description, Operation pseudocode, Flags
Affected, Exceptions); Volume 3 is the \textbf{system programming
guide} covering memory ordering, completion fences, and the host
driver surface.

\RefSection{Notational Conventions}

\begin{itemize}[label={\color{oceanblue}\textbullet}, leftmargin=*, itemsep=0.3ex]
  \item Bit ranges use the \texttt{[MSB:LSB]} convention, inclusive on
        both ends. \texttt{[63:60]} is a 4-bit field.
  \item Hexadecimal constants use the SystemVerilog literal form,
        e.g.\ \texttt{4'h0} (a 4-bit hex zero).
  \item The symbol $\gets$ is used in Operation pseudocode to denote
        assignment; $\equiv$ denotes definitional equivalence.
  \item \texttt{dest} always precedes \texttt{src} in field ordering.
  \item A field annotated ``reserved'' must be written as zero by
        software. The decoder does not currently check this, but future
        revisions may assign the field.
\end{itemize}

\RefSection{Relationship to Other Documentation}

The pccx web portal at \url{https://hwkim-dev.github.io/pccx/}
contains a living, cross-referenced copy of the encoding tables
together with interactive Mermaid, WaveDrom, and Graphviz diagrams.
This PDF is the \textit{offline, signoff-ready} mirror; both sources
are regenerated from the authoritative RTL definitions in the external
repository
\href{https://github.com/hwkim-dev/pccx-FPGA-NPU-LLM-kv260}{\texttt{hwkim-dev/pccx-FPGA-NPU-LLM-kv260}}
under
\texttt{hw/rtl/NPU\_Controller/NPU\_Control\_Unit/ISA\_PACKAGE/isa\_pkg.sv}.
Where this manual and \texttt{isa\_pkg.sv} disagree, the SystemVerilog
package is normative.

\RefSection{Feedback}

File issues against the \texttt{pccx} repository or submit pull
requests to the same. The author reads every report.

\newpage

% ---------- Table of Contents ----------
\tableofcontents
\clearpage
% ---------- List of Figures ----------
\listoffigures
\clearpage

% =================================================================
%  Body
% =================================================================
\pagenumbering{arabic}
\setcounter{chapter}{0}

% =================================================================
\chapter{ABOUT THIS MANUAL}

\section{Scope}

This manual is the authoritative reference for the \textbf{pccx v002
Instruction Set Architecture} (hereafter ``the ISA''). The ISA defines
the exact binary contract between a host driver running on the
Zynq~UltraScale+ Processing System (PS) and the NPU execution complex
that lives on the Programmable Logic (PL). Every opcode, every field,
every reserved bit documented here mirrors
\texttt{isa\_pkg.sv} in the v002 RTL repository.

\section{Audience}

This document assumes familiarity with the following:

\begin{itemize}[label={\color{oceanblue}\textbullet}, leftmargin=*, itemsep=0.3ex]
  \item Systolic arrays with Weight Stationary and Weight Streaming dataflow.
  \item The AXI4 and AXI4-Lite protocols, in particular the Zynq
        UltraScale+ HP and ACP port semantics.
  \item Transformer LLM decoding: Q/K/V projection, scaled
        dot-product attention, the FFN block, RMSNorm, RoPE, and
        token streaming.
  \item Block-floating-point (BFloat16) numerics and INT4 weight
        quantization (WNx/AxAy notation).
\end{itemize}

This manual is \textit{not} a tutorial. For onboarding material,
start from the pccx documentation portal.

\section{Manual Layout}

\begin{center}
\begin{tikzpicture}[
  node distance=0.5cm and 0.55cm,
  vol/.style={rectangle, draw=black, thin,
              fill=oceanblue!8, minimum width=3.5cm, minimum height=1.8cm,
              font=\sffamily\small, align=center, inner sep=4pt},
]
  \node[vol] (v1) {\textbf{Volume 1}\\Basic Architecture\\\textcolor{graytext}{\scriptsize Ch. 2 -- 3}};
  \node[vol, right=of v1] (v2) {\textbf{Volume 2}\\Instruction Set Reference\\\textcolor{graytext}{\scriptsize Ch. 4 -- 6}};
  \node[vol, right=of v2] (v3) {\textbf{Volume 3}\\System Programming Guide\\\textcolor{graytext}{\scriptsize Ch. 7 -- 8}};
  \node[vol, right=of v3] (ap) {\textbf{Appendices}\\A -- D\\\textcolor{graytext}{\scriptsize Opcode Map, etc.}};

  \draw[arr] (v1) -- (v2);
  \draw[arr] (v2) -- (v3);
  \draw[arr] (v3) -- (ap);
\end{tikzpicture}
\captionof{figure}{Manual layout — four volumes bound into one book.}
\label{fig:manual-layout}
\end{center}

\newpage

% =================================================================
\chapter{ARCHITECTURAL OVERVIEW}

This chapter provides the architectural context required to read
subsequent chapters. It describes the host--device boundary, the
internal datapath of the NPU, the memory hierarchy, and the
execution pipeline at a block level. Subsequent chapters refer back
to the figures defined here.

\section{System Context}

The pccx NPU sits inside the PL of the Zynq UltraScale+ and is
driven by the ARM Cortex-A53 cluster on the PS. The host issues
instructions via an AXI4-Lite command FIFO and receives completion
status via a second AXI4-Lite FIFO. Bulk data crosses via four
AXI4 High-Performance ports (HP0--HP3) and one Accelerator Coherency
Port (ACP) for low-latency transfers.

\begin{figure}[h]
\centering
\begin{tikzpicture}[
  node distance = 0.45cm and 1.0cm,
  font = \sffamily\small,
]
  % ---- PS side (uniform size: 3.0 cm × 1.0 cm) ----
  \node[blocknode, minimum width=3.0cm, minimum height=1.0cm, fill=oceanblue!10] (cpu)
    {ARM Cortex-A53\\\scriptsize 4 cores};
  \node[blocknode, below=of cpu, minimum width=3.0cm, minimum height=1.0cm, fill=oceanblue!10]
    (drv) {Host Driver\\\scriptsize libpccx};
  \node[blocknode, below=of drv, minimum width=3.0cm, minimum height=1.0cm, fill=accentpurple!8]
    (ddr) {DDR4\\\scriptsize 4\,GB};
  \node[fit={(cpu)(drv)(ddr)}, draw=black, thin, dashed,
        inner sep=10pt,
        label={[font=\sffamily\bfseries\footnotesize, anchor=south]above:Processing System (PS)}]
        (ps) {};

  % ---- PL side (uniform size: 3.2 cm × 1.0 cm) ----
  \node[blocknode, right=4.0cm of cpu, minimum width=3.2cm, minimum height=1.0cm, fill=accentred!8]
    (ctrl) {Control Unit\\\scriptsize decode · dispatch};
  \node[blocknode, below=of ctrl, minimum width=3.2cm, minimum height=1.0cm, fill=accentpurple!8]
    (comp) {Compute Cluster\\\scriptsize GEMM · GEMV · SFU};
  \node[blocknode, below=of comp, minimum width=3.2cm, minimum height=1.0cm, fill=accentpurple!8]
    (l2) {L2 Cache\\\scriptsize 1.75\,MB URAM};
  \node[fit={(ctrl)(comp)(l2)}, draw=black, thin, dashed,
        inner sep=10pt,
        label={[font=\sffamily\bfseries\footnotesize, anchor=south]above:Programmable Logic (PL) --- pccx NPU}]
        (pl) {};

  % ---- AXI links ----
  % CMD_IN: PS → PL (single straight arrow)
  \draw[arrfat] (cpu.east) --
    node[above, font=\sffamily\scriptsize, fill=white, inner sep=1pt]{AXI4-Lite CMD\_IN \quad 64\,b}
    (ctrl.west);
  % STAT_OUT: PL → PS (return path below CMD_IN)
  \draw[arrfat] (ctrl.west) ++(0,-0.30cm) -- ++(-1.6cm,0)
    node[above, pos=0.5, font=\sffamily\scriptsize, fill=white, inner sep=1pt]{STAT\_OUT}
    -- ++(0,-0.50cm) -| (drv.east);
  % AXI-HP bulk data (DDR ↔ L2)
  \draw[bus, <->] (ddr.east) --
    node[above, font=\sffamily\scriptsize, fill=white, inner sep=1pt]{AXI-HP $\times$4 \enspace 128\,b @ 250\,MHz}
    (l2.west);
\end{tikzpicture}
\caption{Host--device partition. PS issues 64-bit instructions to the
Control Unit over AXI4-Lite, while bulk activations and weights cross
via four 128-bit HP ports.}
\label{fig:system-context}
\end{figure}

\section{NPU Top-Level Block Diagram}

Figure~\ref{fig:npu-top} expands the PL side of Figure~\ref{fig:system-context}.
The Control Unit decodes each 64-bit instruction in one cycle and
emits micro-ops (\textit{\uops}) into the Dispatcher, which resolves
hazards and feeds the GEMM, GEMV, SFU, and DMA back-end resources.

\begin{figure}[h]
\centering
\begin{tikzpicture}[
  node distance = 0.45cm and 0.9cm,
  font = \sffamily\small,
]
  % Row 1 — Control path
  \node[blocknode, minimum width=2.5cm] (cmd) {CMD\_IN\\FIFO};
  \node[blocknode, right=of cmd, minimum width=2.2cm] (dec) {Decoder};
  \node[hotblock,  right=of dec, minimum width=2.6cm] (disp) {Dispatcher};
  \node[scheduler, above=0.35cm of disp] (gs) {Global Scheduler};

  % Row 2 — Back-end resources
  \node[blocknode, below=1.0cm of disp, xshift=-4.8cm, minimum width=2.3cm, fill=accentpurple!8] (gemm) {GEMM\\\scriptsize 32$\times$32 SA};
  \node[blocknode, right=0.35cm of gemm, minimum width=2.3cm, fill=accentpurple!8] (gemv) {GEMV\\\scriptsize 4$\times$32 MAC};
  \node[blocknode, right=0.35cm of gemv, minimum width=2.3cm, fill=accentpurple!8] (sfu)  {SFU (CVO)\\\scriptsize CORDIC/LUT};
  \node[blocknode, right=0.35cm of sfu, minimum width=2.3cm, fill=accentorange!10] (dma) {DMA\\\scriptsize ACP/HP};

  % Row 3 — Memory
  \node[ram, below=1.0cm of gemm, minimum width=2.3cm] (l1g) {L1\\\scriptsize GEMM};
  \node[ram, below=1.0cm of gemv, minimum width=2.3cm] (l1v) {L1\\\scriptsize GEMV};
  \node[ram, below=1.0cm of sfu,  minimum width=2.3cm] (l1s) {L1\\\scriptsize SFU};
  \node[ram, below=1.0cm of dma,  minimum width=2.3cm] (wfif){Weight FIFO\\\scriptsize 4$\times$URAM};

  % L2
  \node[ram, below=0.9cm of l1v, minimum width=10cm, minimum height=1cm, fill=accentpurple!18]
    (l2) {L2 Cache (Unified FMAP + KV) --- 1.75\,MB URAM};
  \node[ram, below=0.4cm of l2, minimum width=10cm, minimum height=0.8cm, fill=skytint]
    (cc) {Constant Cache --- 64 entries};

  % Arrows
  \draw[arr] (cmd) -- (dec);
  \draw[arr] (dec) -- node[above, font=\sffamily\scriptsize]{\uop} (disp);
  \draw[arr, <->] (disp) -- (gs);

  \draw[arr] (disp.south) -- ++(0,-0.35cm) -| (gemm.north);
  \draw[arr] (disp.south) -- ++(0,-0.35cm) -| (gemv.north);
  \draw[arr] (disp.south) -- ++(0,-0.35cm) -| (sfu.north);
  \draw[arr] (disp.south) -- ++(0,-0.35cm) -| (dma.north);

  \draw[bus] (gemm) -- (l1g);
  \draw[bus] (gemv) -- (l1v);
  \draw[bus] (sfu)  -- (l1s);
  \draw[bus] (dma)  -- (wfif);

  \draw[bus] (l1g) -- (l2);
  \draw[bus] (l1v) -- (l2);
  \draw[bus] (l1s) -- (l2);
  \draw[bus] (wfif) -- (l2);
  \draw[arr] (cc.east) -| ++(0.6cm,1.6cm) |- (disp.east);
\end{tikzpicture}
\caption{NPU top-level block diagram. Control flows top-down; data
flows bottom-up via the unified L2 cache.}
\label{fig:npu-top}
\end{figure}

\section{Clock Domains and CDC}

Two clock domains meet inside the NPU: the 250\,MHz AXI domain
adjacent to the PS, and the 400\,MHz compute core. A bank of Xilinx
Parameterized Macro (XPM) asynchronous FIFOs in
\texttt{independent\_clock} mode is the only path that crosses the
boundary; there is no manual synchronizer logic.

\begin{figure}[h]
\centering
\begin{tikzpicture}[
  node distance = 0.5cm and 0.9cm,
  font = \sffamily\small,
]
  \node[blocknode, minimum width=2.6cm, minimum height=0.9cm, fill=accentorange!10]
    (wr) {Write\\\scriptsize 250\,MHz};
  \node[blocknode, right=1.0cm of wr, minimum width=3.6cm, minimum height=0.9cm, fill=black!4]
    (fifo) {XPM Async FIFO\\\scriptsize 4096$\times$128\,b};
  \node[blocknode, right=1.0cm of fifo, minimum width=2.6cm, minimum height=0.9cm, fill=accentgreen!10]
    (rd) {Read\\\scriptsize 400\,MHz};

  % Arrows — labels with white fill to prevent overlap
  \draw[arrfat] (wr) -- node[above, font=\sffamily\scriptsize, fill=white, inner sep=1.5pt]{data + wr\_en} (fifo.west);
  \draw[arrfat] (fifo.east) -- node[above, font=\sffamily\scriptsize, fill=white, inner sep=1.5pt]{data + rd\_en} (rd);

  % Clock labels — placed below each block to avoid overlapping arrow labels
  \node[font=\sffamily\scriptsize, below=0.18cm of wr]  {CLK\_PS (250\,MHz)};
  \node[font=\sffamily\scriptsize, below=0.18cm of rd]  {CLK\_PL (400\,MHz)};

  % CDC boundary dashed lines (at FIFO edges)
  \draw[black!50, dashed] ($(fifo.north west)+(0,0.35cm)$) -- ($(fifo.south west)+(0,-0.55cm)$);
  \draw[black!50, dashed] ($(fifo.north east)+(0,0.35cm)$) -- ($(fifo.south east)+(0,-0.55cm)$);

  \node[domainlabel, above=0.45cm of wr]   {PS clock domain};
  \node[domainlabel, above=0.45cm of rd]   {PL clock domain};
\end{tikzpicture}
\caption{Clock-domain crossing path: a 4096$\times$128\,b XPM FIFO in
\texttt{independent\_clock} mode decouples the 250\,MHz AXI side from
the 400\,MHz compute core.}
\label{fig:cdc}
\end{figure}

\subsection*{AXI-HP Read Handshake}

Every weight or activation byte that crosses from the PS side of the
chip to the PL side does so over the AXI-HP 4-channel burst-capable
read protocol.  The ingress direction of one such handshake (single
beat shown; a full burst chains 16 of these on the data phase) is
illustrated in Figure~\ref{fig:axi-timing}.  The protocol enforces
three rules the software driver must respect:

\begin{itemize}[label={\color{oceanblue}\textbullet}, leftmargin=*, itemsep=0.3ex]
    \item \texttt{ARVALID} must remain asserted until both
          \texttt{ARVALID} and \texttt{ARREADY} are sampled high on
          the same rising edge (the \textit{address handshake}).
    \item \texttt{ARADDR} must be stable across that window.
    \item After the handshake, the master must be ready to accept
          \texttt{RDATA} within at most \texttt{T\_rdata\_max} cycles
          of the corresponding \texttt{RVALID} assertion.
\end{itemize}

\begin{figure}[h]
\centering
\begin{tikzpicture}[
  font=\sffamily\scriptsize,
  sig/.style={semithick, draw=oceanblue, line cap=round},
  clk/.style={semithick, draw=graytext, line cap=round},
  grid/.style={very thin, dotted, draw=graytext!55},
  hilite/.style={thick, draw=accentred, dashed},
  lbl/.style={anchor=east, font=\sffamily\small},
]
  \def\cw{0.9}   % cycle width

  % Signal labels (left column)
  \node[lbl] at (-0.1,  5) {CLK};
  \node[lbl] at (-0.1,  4) {ARVALID};
  \node[lbl] at (-0.1,  3) {ARREADY};
  \node[lbl] at (-0.1,  2) {ARADDR[16:0]};
  \node[lbl] at (-0.1,  1) {RVALID};
  \node[lbl] at (-0.1,  0) {RDATA[127:0]};

  % Cycle grid + numbers
  \foreach \i in {0,...,6} {
    \draw[grid] ({\cw*\i}, -0.4) -- ({\cw*\i}, 5.6);
    \node[font=\sffamily\scriptsize\color{graytext}] at ({\cw*\i + \cw/2}, 5.75) {t\the\numexpr\i+1\relax};
  }

  % Clock (6 cycles), square wave with half-period of 0.5 cycle
  \draw[clk] (0, 5.0)
    -- ({\cw*0.5}, 5.0) -- ({\cw*0.5}, 5.4) -- ({\cw*1.0}, 5.4) -- ({\cw*1.0}, 5.0)
    -- ({\cw*1.5}, 5.0) -- ({\cw*1.5}, 5.4) -- ({\cw*2.0}, 5.4) -- ({\cw*2.0}, 5.0)
    -- ({\cw*2.5}, 5.0) -- ({\cw*2.5}, 5.4) -- ({\cw*3.0}, 5.4) -- ({\cw*3.0}, 5.0)
    -- ({\cw*3.5}, 5.0) -- ({\cw*3.5}, 5.4) -- ({\cw*4.0}, 5.4) -- ({\cw*4.0}, 5.0)
    -- ({\cw*4.5}, 5.0) -- ({\cw*4.5}, 5.4) -- ({\cw*5.0}, 5.4) -- ({\cw*5.0}, 5.0)
    -- ({\cw*5.5}, 5.0) -- ({\cw*5.5}, 5.4) -- ({\cw*6.0}, 5.4) -- ({\cw*6.0}, 5.0);

  % ARVALID: asserted cycles t2..t3 (master), deassertion on t4
  \draw[sig] (0, 4) -- ({\cw*1}, 4)
    -- ({\cw*1}, 4.4) -- ({\cw*3}, 4.4) -- ({\cw*3}, 4)
    -- ({\cw*6}, 4);

  % ARREADY: asserted cycle t3 (slave), 1-cycle high
  \draw[sig] (0, 3) -- ({\cw*2}, 3)
    -- ({\cw*2}, 3.4) -- ({\cw*3}, 3.4) -- ({\cw*3}, 3)
    -- ({\cw*6}, 3);

  % ARADDR: valid t2..t3 with addr label
  \draw[sig] (0, 2) -- ({\cw*1}, 2);
  \draw[sig] ({\cw*1}, 2) -- ({\cw*1.1}, 2.4) -- ({\cw*3}, 2.4)
             -- ({\cw*3.1}, 2) -- ({\cw*6}, 2);
  \node[font=\sffamily\scriptsize] at ({\cw*2}, 2.2) {\texttt{addr}};

  % RVALID: asserted t5 (slave returns data one cycle after handshake+latency)
  \draw[sig] (0, 1) -- ({\cw*4}, 1)
    -- ({\cw*4}, 1.4) -- ({\cw*5}, 1.4) -- ({\cw*5}, 1)
    -- ({\cw*6}, 1);

  % RDATA: valid during RVALID
  \draw[sig] (0, 0) -- ({\cw*4}, 0);
  \draw[sig] ({\cw*4}, 0) -- ({\cw*4.1}, 0.4) -- ({\cw*5}, 0.4)
             -- ({\cw*5.1}, 0) -- ({\cw*6}, 0);
  \node[font=\sffamily\scriptsize] at ({\cw*4.5}, 0.2) {\texttt{data}};

  % Handshake marker at t3 (both ARVALID & ARREADY high)
  \draw[hilite] ({\cw*2.5}, 4.55) -- ({\cw*2.5}, 2.85);
  \node[font=\sffamily\scriptsize\color{accentred}, anchor=west]
    at ({\cw*2.6}, 3.7) {address handshake};

  % Data-phase marker at t5
  \draw[hilite] ({\cw*4.5}, 1.55) -- ({\cw*4.5}, -0.25);
  \node[font=\sffamily\scriptsize\color{accentred}, anchor=west]
    at ({\cw*4.6}, 0.7) {data beat};
\end{tikzpicture}
\caption{AXI-HP read protocol, single beat.  On cycle \texttt{t3} both
\texttt{ARVALID} and \texttt{ARREADY} are high --- the address is
captured.  Two cycles of round-trip latency later (cycle \texttt{t5})
the slave returns \texttt{RDATA} with \texttt{RVALID}.  A full burst
chains 16 data beats on back-to-back cycles once the address phase
has committed.}
\label{fig:axi-timing}
\end{figure}

\section{Instruction Pipeline}

\begin{figure}[h]
\centering
\begin{tikzpicture}[
  node distance = 0.25cm,
  font = \sffamily\footnotesize,
]
  \node[stage] (fe)  {IF\\\scriptsize Fetch};
  \node[stage, right=of fe]  (de) {ID\\\scriptsize Decode};
  \node[stage, right=of de]  (ds) {DS\\\scriptsize Dispatch};
  \node[stage, right=of ds, fill=accentpurple!15, draw=accentpurple]  (ex) {EX\\\scriptsize Execute};
  \node[stage, right=of ex]  (mem){ME\\\scriptsize Memory};
  \node[stage, right=of mem] (wb) {WB\\\scriptsize Writeback};
  \node[stage, right=of wb, fill=accentgreen!15, draw=accentgreen]  (co) {CO\\\scriptsize Commit};

  \foreach \a/\b in {fe/de, de/ds, ds/ex, ex/mem, mem/wb, wb/co}
    \draw[arr] (\a) -- (\b);

  % Hazard feedback
  \draw[ctrl] (ds.north) ++(0,0.1cm) -- ++(0,0.5cm)
    -| node[pos=0.25, above, font=\sffamily\scriptsize]{stall on hazard} (fe.north);
\end{tikzpicture}
\caption{Conceptual 7-stage NPU pipeline. Because each instruction can
launch thousands of MAC cycles, the \textbf{Execute} stage dominates
real latency; the front end merely streams instructions in as fast as
resources free up.}
\label{fig:pipeline}
\end{figure}

\section{Memory Hierarchy}

\begin{figure}[h]
\centering
\begin{tikzpicture}[
  node distance = 0.3cm,
  font = \sffamily\small,
]
  % Layer stack
  \node[rectangle, draw=black, thin, fill=black!8,         minimum width=10cm, minimum height=0.9cm] (ddr) {DDR4 --- 4\,GB @ 12.8\,GB/s};
  \node[rectangle, draw=black, thin, fill=accentorange!10, minimum width=8.5cm, minimum height=0.9cm, below=of ddr] (acp) {AXI-HP $\times$4 + ACP --- 16\,GB/s burst};
  \node[rectangle, draw=black, thin, fill=accentpurple!10, minimum width=7.0cm, minimum height=0.9cm, below=of acp] (l2) {Unified L2 --- 1.75\,MB URAM (FMAP + KV)};
  \node[rectangle, draw=black, thin, fill=oceanblue!10,    minimum width=5.5cm, minimum height=0.9cm, below=of l2]  (l1) {Per-core L1 --- 128\,KB BRAM};
  \node[rectangle, draw=black, thin, fill=accentgreen!10,  minimum width=4.0cm, minimum height=0.9cm, below=of l1]  (rf) {Systolic Array / RegFile};

  % Latency labels (right side)
  \node[font=\sffamily\scriptsize\color{graytext}, right=0.2cm of ddr] {$\sim$200 cyc};
  \node[font=\sffamily\scriptsize\color{graytext}, right=0.2cm of acp] {$\sim$40 cyc};
  \node[font=\sffamily\scriptsize\color{graytext}, right=0.2cm of l2]  {$\sim$6 cyc};
  \node[font=\sffamily\scriptsize\color{graytext}, right=0.2cm of l1]  {$\sim$2 cyc};
  \node[font=\sffamily\scriptsize\color{graytext}, right=0.2cm of rf]  {0 cyc};

  % Capacity labels (left)
  \node[font=\sffamily\scriptsize\color{graytext}, left=0.2cm of ddr] {$2^{32}$ B};
  \node[font=\sffamily\scriptsize\color{graytext}, left=0.2cm of l2]  {$2^{21}$ B};
  \node[font=\sffamily\scriptsize\color{graytext}, left=0.2cm of l1]  {$2^{17}$ B};
  \node[font=\sffamily\scriptsize\color{graytext}, left=0.2cm of rf]  {$2^{10}$ B};
\end{tikzpicture}
\caption{Memory hierarchy pyramid. Latency increases upward; capacity
increases upward; bandwidth decreases upward. The unified 1.75\,MB L2
dominates the on-chip budget.}
\label{fig:mem-hier}
\end{figure}

\begin{SideNote}{URAM allocation breakdown}
Of the 64 URAM blocks on the KV260, 8 are claimed by the 4-port
Weight FIFO (2 per port) and the remaining 56 are paired into 28
structural pairs to form the unified L2 (28 pairs $\times$ 4096
depth $\times$ 128-bit = 114\,688 entries $\times$ 16 B =
$\sim$1.75 MB). See Figure~\ref{fig:uram-alloc}.
\end{SideNote}

\begin{figure}[h]
\centering
\begin{tikzpicture}[
  font = \sffamily\scriptsize,
]
  % 64 cells in an 8×8 grid (cell size 0.6 cm, gap 0.05 cm → step 0.65 cm)
  \foreach \i in {0,...,7} {
    \foreach \j in {0,...,7} {
      \pgfmathtruncatemacro{\n}{\i*8+\j}
      \pgfmathtruncatemacro{\isWeight}{\n < 8 ? 1 : 0}
      \ifnum\isWeight=1
        \draw[fill=accentorange!30, draw=black, thin] (\j*0.65,-\i*0.65) rectangle ++(0.58,0.58);
      \else
        \draw[fill=accentpurple!20, draw=black, thin] (\j*0.65,-\i*0.65) rectangle ++(0.58,0.58);
      \fi
      \node at (\j*0.65 + 0.29, -\i*0.65 + 0.29) {\n};
    }
  }

  % Legend — placed below the grid (grid bottom ≈ y = -(7*0.65+0.58) = -5.13)
  \draw[fill=accentorange!30, draw=black, thin] (0,-5.5) rectangle ++(0.45,0.45);
  \node[anchor=west, font=\sffamily\small] at (0.6,-5.27)
    {Weight FIFO (8 URAMs, 2 per port $\times$ 4 ports)};
  \draw[fill=accentpurple!20, draw=black, thin] (0,-6.15) rectangle ++(0.45,0.45);
  \node[anchor=west, font=\sffamily\small] at (0.6,-5.92)
    {Unified L2 --- FMAP + KV Cache (56 URAMs, 28 pairs)};
\end{tikzpicture}
\caption{Physical URAM allocation on the Kria KV260 (64 blocks total).
Blocks 0--7 form the 4-port Weight FIFO; blocks 8--63 pair up to
create the 1.75 MB unified L2.}
\label{fig:uram-alloc}
\end{figure}

\newpage

% =================================================================
\chapter{NUMERICS AND EXECUTION MODEL}

\section{Block Floating Point + INT4 Weight Quantization}

pccx performs dot products with exact integer arithmetic by aligning
activation exponents to a common Block Floating Point (BF16) scale
factor, then multiplying by INT4-quantized weights. A single
DSP48E2 executes the 9-bit $\times$ 4-bit inner multiplication
natively; the BF16 scale collapses out to a cheap scalar multiply
at the end of the accumulation.

For two activation--weight pairs sharing scale factors
$S_{\mathrm{fmap}}$ (for BF16) and $S_{\mathrm{weight}}$ (for INT4):

\begin{equation}
\begin{aligned}
A &= a \cdot S_{\mathrm{fmap}}, \quad  &  B &= b \cdot S_{\mathrm{fmap}} \\
C &= c \cdot S_{\mathrm{weight}}, \quad  &  D &= d \cdot S_{\mathrm{weight}}
\end{aligned}
\end{equation}

The quantity $(A \cdot C) + (B \cdot D)$ factors to:

\begin{equation}
\mathrm{Result} = \underbrace{(a \cdot c + b \cdot d)}_{\text{DSP48E2 integer MAC}} \cdot \underbrace{S_{\mathrm{fmap}} \cdot S_{\mathrm{weight}}}_{\text{post-process scalar}}
\end{equation}

\begin{figure}[h]
\centering
\begin{tikzpicture}[
  font = \sffamily\small,
  node distance = 0.55cm and 1.0cm,
]
  % Inputs
  \node[blocknode, fill=accentorange!10] (act) {Activation\\\scriptsize BF16};
  \node[blocknode, below=of act, fill=accentpurple!8] (w)  {Weight\\\scriptsize INT4};

  % Intermediate stages — extra horizontal gap (1.1 cm) keeps them clear of DSP
  \node[blocknode, right=1.1cm of act] (emax)  {$E_{\max}$\\\scriptsize align};
  \node[blocknode, right=1.1cm of w]   (signx) {Sign-extend\\\scriptsize INT4$\to$INT9};

  % DSP48E2 — centred between emax and signx, with extra gap (1.4 cm)
  \node[blocknode, fill=accentred!8,
        right=1.4cm of $(emax)!0.5!(signx)$,
        minimum height=2.3cm, minimum width=2.5cm] (dsp)
    {DSP48E2\\\scriptsize $A \cdot B + C$\\\scriptsize 48-bit ACC};

  % Post-process and output
  \node[blocknode, right=1.0cm of dsp, minimum width=2.8cm] (post)
    {Post-process\\$\times(S_f \cdot S_w)$};
  \node[blocknode, right=1.0cm of post, fill=accentgreen!10] (o) {Output\\\scriptsize BF16};

  \draw[arr] (act)   -- (emax);
  \draw[arr] (w)     -- (signx);
  \draw[arr] (emax)  -- (dsp.west |- emax);
  \draw[arr] (signx) -- (dsp.west |- signx);
  \draw[arr] (dsp)   -- (post);
  \draw[arr] (post)  -- (o);
\end{tikzpicture}
\caption{DSP48E2 datapath. INT9 $\times$ INT4 multiplication executes
natively; the BF16 scale factor product is applied once per MAC epoch
by the post-process stage.}
\label{fig:dsp-datapath}
\end{figure}

\section{Systolic Dataflow — Weight Stationary (GEMM)}

\begin{figure}[h]
\centering
\begin{tikzpicture}[
  font = \sffamily\scriptsize,
  % PE cells: 0.70 cm square, 0.80 cm step → 0.10 cm clear gap between cells
  pe/.style={rectangle, draw=black, thin, fill=oceanblue!8,
             minimum width=0.70cm, minimum height=0.70cm, inner sep=0pt},
]
  \def\step{0.80}   % grid step (cm)
  \def\half{0.35}   % half PE size

  % ---- 6×6 PE array — weights labelled inside each PE ----
  \foreach \i in {0,...,5}
    \foreach \j in {0,...,5} {
      \node[pe] (pe\i\j) at (\j*\step, -\i*\step) {$W_{\i\j}$};
    }

  % ---- Activation inputs from left ----
  \foreach \i in {0,...,5} {
    \draw[arr] ({-1.0}, {-\i*\step}) -- (pe\i0.west);
    \node[anchor=east, font=\scriptsize] at (-1.05, -\i*\step) {$a_\i$};
  }

  % ---- Partial-sum inputs from top (arrows end at top edge of row-0 PEs) ----
  \foreach \j in {0,...,5} {
    \draw[arr] ({\j*\step}, {0.70}) -- ({\j*\step}, {\half+0.02});
    \node[anchor=south, font=\scriptsize] at (\j*\step, 0.72) {$\psi_\j^{\mathrm{in}}$};
  }

  % ---- Partial-sum outputs from bottom (arrows start at bottom edge of row-5 PEs) ----
  \def\botY{{-5*\step}}
  \foreach \j in {0,...,5} {
    \draw[arr] ({\j*\step}, {-5*\step-\half-0.02}) -- ({\j*\step}, {-5*\step-0.72});
    \node[anchor=north, font=\scriptsize] at (\j*\step, {-5*\step-0.74}) {$\psi_\j^{\mathrm{out}}$};
  }

  % ---- Activation pass-through to right ----
  \foreach \i in {0,...,5} {
    \draw[arr] (pe\i5.east) -- ++({0.55},0);
  }

  % ---- Annotation ----
  \node[anchor=north west, font=\sffamily\small, align=left]
    at (0, {-5*\step-1.0}) {Weights $W_{ij}$ preloaded per tile (Weight Stationary).\\
                             Activations $a_i$ stream left-to-right; partial sums $\psi$ top-to-bottom.};
\end{tikzpicture}
\caption{Systolic array dataflow (illustrated on a 6$\times$6; actual
hardware is 32$\times$32). Weights are preloaded once per tile;
activations and partial sums march through the array every cycle.}
\label{fig:systolic}
\end{figure}

\section{Systolic Dataflow — Weight Streaming (GEMV)}

GEMV differs from GEMM in that the weight matrix is streamed through
4 parallel cores row-partition-wise rather than pinned inside PEs. A
single GEMV dispatches 32-lane LUT-based MAC units followed by a
5-stage reduction tree, producing one scalar per decoded token.

\begin{figure}[h]
\centering
\begin{tikzpicture}[
  font = \sffamily\small,
  node distance = 0.35cm and 0.6cm,
  % All four cores: identical width + height for visual uniformity
  core/.style={rectangle, draw=black, thin,
               fill=accentpurple!8,
               minimum width=2.6cm, minimum height=1.1cm,
               font=\sffamily\small, align=center},
]
  \node[blocknode, minimum width=2.0cm, minimum height=1.1cm, fill=accentpurple!8]
    (wbuf) {Weight FIFO\\\scriptsize URAM};

  % Four cores at fixed y offsets — identical node style
  \node[core, right=1.4cm of wbuf, yshift= 1.65cm] (c0) {GEMV Core 0};
  \node[core, right=1.4cm of wbuf, yshift= 0.55cm] (c1) {GEMV Core 1};
  \node[core, right=1.4cm of wbuf, yshift=-0.55cm] (c2) {GEMV Core 2};
  \node[core, right=1.4cm of wbuf, yshift=-1.65cm] (c3) {GEMV Core 3};

  % Reduction / concat stage — centred between c1 and c2
  \node[blocknode, right=1.2cm of c1, yshift=-0.55cm,
        minimum width=2.5cm, minimum height=1.1cm, fill=accentgreen!10]
    (red) {Row-partition\\concat};
  \node[blocknode, right=0.8cm of red, minimum width=1.8cm, minimum height=1.1cm,
        fill=accentpurple!8] (l2) {L2\\result};

  % Weight bus (trunk style: horizontal line from wbuf, then vertical drops)
  \coordinate (wjct) at ($(wbuf.east)+(0.5cm,0)$);
  \draw[-] (wbuf.east) -- (wjct);
  \draw[trunk] ($(wjct)+(0, 1.65cm)$) -- ($(wjct)+(0,-1.65cm)$);
  \foreach \c in {c0,c1,c2,c3} \draw[arrfat] (wjct |- \c) -- (\c.west);

  % Core → reduce arrows
  \foreach \c in {c0,c1,c2,c3} \draw[arr] (\c.east) -- (red.west |- \c);

  \draw[arr] (red) -- (l2);
\end{tikzpicture}
\caption{GEMV weight streaming: the 4 GEMV cores each consume a
row-partition stream from the URAM Weight FIFO; results concatenate
into a single output vector in L2.}
\label{fig:gemv-stream}
\end{figure}

Each GEMV core reduces its 32 partial products to a single scalar
through a five-stage binary tree (Figure~\ref{fig:gemv-tree}).  Stage
1 is implemented with 16 DSP48E2 slices running in $A+B$ mode;
stages 2--5 fall to 6-LUT adders to conserve the DSP budget for the
GEMM array.

\begin{figure}[h]
\centering
\resizebox{\textwidth}{!}{%
\begin{tikzpicture}[
  font=\sffamily\scriptsize,
  leaf/.style={rectangle, draw=black, thin, fill=black!6,
               minimum width=0.42cm, minimum height=0.42cm, inner sep=0pt},
  dsp/.style ={rectangle, draw=black, thin, fill=accentred!10,
               minimum width=0.6cm, minimum height=0.5cm, inner sep=0pt,
               font=\sffamily\tiny\bfseries},
  lut/.style ={rectangle, draw=black, thin, fill=accentgreen!10,
               minimum width=0.75cm, minimum height=0.5cm, inner sep=0pt,
               font=\sffamily\tiny\bfseries},
  sink/.style={rectangle, draw=black, thin, fill=accentpurple!12,
               minimum width=0.9cm, minimum height=0.55cm, inner sep=0pt,
               font=\sffamily\tiny\bfseries},
  rowlbl/.style={anchor=east, font=\sffamily\scriptsize\color{graytext}},
]
  % Stage 0: 32 leaves (partial products), spaced 0.5cm apart
  \def\leafdx{0.5}
  \foreach \i in {0,...,31}
    \node[leaf] (p\i) at (\i*\leafdx, 0) {};
  \node[rowlbl] at (-0.3, 0)   {32 $\times$ W$\cdot$A};

  % Stage 1 (DSP48E2): 16 nodes, each merging 2 leaves -> positioned midway
  \foreach \i in {0,...,15} {
    \pgfmathtruncatemacro{\a}{2*\i}
    \pgfmathtruncatemacro{\b}{2*\i+1}
    \pgfmathsetmacro{\xc}{(\a+\b)*\leafdx/2}
    \node[dsp] (d\i) at (\xc, -1.0) {D};
    \draw[arr, -] (p\a) -- (d\i);
    \draw[arr, -] (p\b) -- (d\i);
  }
  \node[rowlbl] at (-0.3, -1.0) {Stage 1 (DSP48E2) 32 $\to$ 16};

  % Stage 2 (LUT): 8 nodes
  \foreach \i in {0,...,7} {
    \pgfmathtruncatemacro{\a}{2*\i}
    \pgfmathtruncatemacro{\b}{2*\i+1}
    \pgfmathsetmacro{\xc}{(\a+\b)*2*\leafdx/2 + \leafdx/2}
    \node[lut] (e\i) at (\xc, -2.0) {L};
    \draw[arr, -] (d\a) -- (e\i);
    \draw[arr, -] (d\b) -- (e\i);
  }
  \node[rowlbl] at (-0.3, -2.0) {Stage 2 (LUT) 16 $\to$ 8};

  % Stage 3 (LUT): 4 nodes
  \foreach \i in {0,...,3} {
    \pgfmathtruncatemacro{\a}{2*\i}
    \pgfmathtruncatemacro{\b}{2*\i+1}
    \pgfmathsetmacro{\xc}{(\a+\b)*4*\leafdx/2 + 1.5*\leafdx}
    \node[lut] (f\i) at (\xc, -3.0) {L};
    \draw[arr, -] (e\a) -- (f\i);
    \draw[arr, -] (e\b) -- (f\i);
  }
  \node[rowlbl] at (-0.3, -3.0) {Stage 3 (LUT) 8 $\to$ 4};

  % Stage 4 (LUT): 2 nodes
  \foreach \i in {0,1} {
    \pgfmathtruncatemacro{\a}{2*\i}
    \pgfmathtruncatemacro{\b}{2*\i+1}
    \pgfmathsetmacro{\xc}{(\a+\b)*8*\leafdx/2 + 3.5*\leafdx}
    \node[lut] (g\i) at (\xc, -4.0) {L};
    \draw[arr, -] (f\a) -- (g\i);
    \draw[arr, -] (f\b) -- (g\i);
  }
  \node[rowlbl] at (-0.3, -4.0) {Stage 4 (LUT) 4 $\to$ 2};

  % Stage 5 (LUT): 1 scalar
  \node[sink] (s) at (31*\leafdx/2, -5.0) {scalar};
  \draw[arr, -] (g0) -- (s);
  \draw[arr, -] (g1) -- (s);
  \node[rowlbl] at (-0.3, -5.0) {Stage 5 (LUT) 2 $\to$ 1};
\end{tikzpicture}%
}
\caption{Per-core GEMV reduction tree (one of four cores shown).
Stage 1 absorbs 32 partial products into 16 sums using DSP48E2 $A+B$
adders; stages 2--5 use 6-LUT adders.  Each stage carries a pipeline
register so the tree closes timing at 400\,MHz.}
\label{fig:gemv-tree}
\end{figure}

\newpage

% =================================================================
\chapter{INSTRUCTION ENCODING}

\section{Top-Level 64-Bit Layout}

Every pccx v002 instruction is exactly \textbf{64 bits wide} and
encodes a 4-bit opcode in the most-significant nibble followed by a
60-bit instruction body whose shape depends on the opcode.

\begin{figure}[h]
\centering
\begin{bytefield}[bitwidth=0.68em, bitheight=2.2em, boxformatting={\centering\sffamily\scriptsize}]{64}
  \bitheader[endianness=big,lsb=0]{0,15,20,26,43,59,63} \\
  \bitbox{4}{\color{oceanblue}\textbf{opcode}\\[0.2em]\scriptsize[63:60]} &
  \bitbox{60}{\color{graytext}\textbf{instruction body --- opcode-specific}\\[0.2em]\scriptsize[59:0]}
\end{bytefield}
\caption{Top-level 64-bit instruction word.}
\label{fig:top-layout}
\end{figure}

\section{Opcode Map}

Of the 16 possible opcodes, five are defined in v002; the remaining
eleven are reserved. Decoder behaviour on reserved opcodes is
\textbf{undefined}.

\begin{longtable}{|c|l|l|p{7.0cm}|}
\hline
\rowcolor{skytint}
\textbf{Opcode} & \textbf{Mnemonic} & \textbf{Family} & \textbf{Function} \\
\hline
\endfirsthead

\texttt{4'h0} & \texttt{GEMV}   & Type A & Matrix $\times$ vector. Drives 4-core 32-MAC streaming engine. \\
\hline
\texttt{4'h1} & \texttt{GEMM}   & Type A & Matrix $\times$ matrix. Drives 32$\times$32 weight-stationary systolic array. \\
\hline
\texttt{4'h2} & \texttt{MEMCPY} & Type B & Bulk data movement: host $\leftrightarrow$ device, L2 $\leftrightarrow$ L2. \\
\hline
\texttt{4'h3} & \texttt{MEMSET} & Type C & Constant Cache writes (shape / size / scale). \\
\hline
\texttt{4'h4} & \texttt{CVO}    & Type D & Complex Vector Op. Transcendentals + reductions via SFU. \\
\hline
\texttt{4'h5}--\texttt{4'hF} & --- & --- & \textit{Reserved.} \\
\hline
\end{longtable}

\section{Decode and Dispatch Pathway}

\begin{figure}[h]
\centering
\begin{tikzpicture}[
  font = \sffamily\footnotesize,
  node distance = 0.4cm and 0.9cm,
  branch/.style={rectangle, draw=black, thin, fill=oceanblue!8,
                 minimum width=2.4cm, minimum height=0.72cm, align=center,
                 font=\sffamily\footnotesize},
]
  \node[blocknode, minimum width=2.4cm] (cmd) {CMD\_IN FIFO};
  \node[blocknode, right=of cmd, minimum width=2.4cm] (dec) {opcode[63:60]\\branch};

  % 5 branch nodes — same style, same width
  \node[branch, right=1.8cm of dec, yshift= 2.16cm] (b0) {GEMV \uop};
  \node[branch, right=1.8cm of dec, yshift= 1.08cm] (b1) {GEMM \uop};
  \node[branch, right=1.8cm of dec, yshift= 0.00cm] (b2) {MEMCTRL \uop};
  \node[branch, right=1.8cm of dec, yshift=-1.08cm] (b3) {MEMSET \uop};
  \node[branch, right=1.8cm of dec, yshift=-2.16cm] (b4) {CVO \uop};
  \node[font=\sffamily\scriptsize, right=1.8cm of dec, yshift=-3.00cm]
    (res) {reserved $\to$ UB};

  \node[blocknode, fill=accentred!8, right=1.6cm of b2, minimum width=2.4cm]
    (disp) {Dispatcher};
  \node[scheduler, above=0.4cm of disp] (gs) {Hazard check};

  % CMD → dec
  \draw[arr] (cmd) -- (dec);

  % Trunk fan-out: stub → vertical trunk → horizontal branches
  \coordinate (jct) at ($(dec.east)+(0.65cm,0)$);
  \draw[-] (dec.east) -- (jct);
  % Vertical trunk spanning all 5 branches
  \draw[trunk] ($(jct)+(0, 2.16cm)$) -- ($(jct)+(0,-2.16cm)$);
  % Horizontal branches with labels
  \draw[arr] ($(jct)+(0, 2.16cm)$) --
    node[above, font=\sffamily\scriptsize, fill=white, inner sep=1pt]{\texttt{4'h0}} (b0.west);
  \draw[arr] ($(jct)+(0, 1.08cm)$) --
    node[above, font=\sffamily\scriptsize, fill=white, inner sep=1pt]{\texttt{4'h1}} (b1.west);
  \draw[arr] ($(jct)+(0, 0.00cm)$) --
    node[above, font=\sffamily\scriptsize, fill=white, inner sep=1pt]{\texttt{4'h2}} (b2.west);
  \draw[arr] ($(jct)+(0,-1.08cm)$) --
    node[above, font=\sffamily\scriptsize, fill=white, inner sep=1pt]{\texttt{4'h3}} (b3.west);
  \draw[arr] ($(jct)+(0,-2.16cm)$) --
    node[above, font=\sffamily\scriptsize, fill=white, inner sep=1pt]{\texttt{4'h4}} (b4.west);

  % Branches → Dispatcher
  \foreach \b in {b0,b1,b2,b3,b4} \draw[arr] (\b.east) -- (disp.west |- \b);

  \draw[arr, <->] (disp) -- (gs);
\end{tikzpicture}
\caption{Decode and dispatch pathway. The opcode selects one of five
\uop types; the dispatcher resolves inter-instruction hazards before
handing off to the back end.}
\label{fig:decode-dispatch}
\end{figure}

\section{\uop Decomposition Summary}

\begin{longtable}{|l|p{10cm}|}
\hline
\rowcolor{skytint}
\textbf{\uop type} & \textbf{Constituent fields (width)} \\
\hline
\endfirsthead

\texttt{gemm\_control\_uop\_t}   & \texttt{flags}~(6\,b) + \texttt{size\_ptr\_addr}~(6\,b) + \texttt{parallel\_lane}~(5\,b) \\
\hline
\texttt{GEMV\_control\_uop\_t}   & Identical layout to \texttt{gemm\_control\_uop\_t}. \\
\hline
\texttt{memory\_control\_uop\_t} & \texttt{data\_route}~(8\,b) + \texttt{dest\_addr}~(17\,b) + \texttt{src\_addr}~(17\,b) + \texttt{shape\_ptr}~(6\,b) + \texttt{async}~(1\,b) \\
\hline
\texttt{memory\_set\_uop\_t}     & \texttt{dest\_cache}~(2\,b) + \texttt{dest\_addr}~(6\,b) + \texttt{a\_value}~(16\,b) + \texttt{b\_value}~(16\,b) + \texttt{c\_value}~(16\,b) \\
\hline
\texttt{cvo\_control\_uop\_t}    & \texttt{cvo\_func}~(4\,b) + \texttt{src\_addr}~(17\,b) + \texttt{dst\_addr}~(17\,b) + \texttt{length}~(16\,b) + \texttt{flags}~(5\,b) + \texttt{async}~(1\,b) \\
\hline
\end{longtable}

\section{Instruction Format Families — Bit-Level Diagrams}

The 60-bit body falls into four structural families. Figures
\ref{fig:type-a}, \ref{fig:type-b}, \ref{fig:type-c}, and
\ref{fig:type-d} show each layout as a true bit-field diagram.

\subsection{Type A --- Compute Format (GEMV / GEMM)}

\begin{figure}[htb]
\centering
\begin{bytefield}[bitwidth=0.68em, bitheight=2.4em,
                  boxformatting={\centering\sffamily\scriptsize\bfseries}]{64}
  \bitheader[endianness=big,lsb=0]{0,3,8,14,20,26,43,60,63} \\
  \bitbox{4}{opcode\\\scriptsize[63:60]} &
  \bitbox{17}{dest\_reg\\\scriptsize[59:43]} &
  \bitbox{17}{src\_addr\\\scriptsize[42:26]} &
  \bitbox{6}{flags\\\scriptsize[25:20]} &
  \bitbox{6}{size\_ptr\\\scriptsize[19:14]} &
  \bitbox{6}{shape\_ptr\\\scriptsize[13:8]} &
  \bitbox{5}{lane\\\scriptsize[7:3]} &
  \bitbox{3}{rsv\\\scriptsize[2:0]}
\end{bytefield}
\caption{\textbf{Type A} body layout, used by \texttt{GEMV} (opcode
\texttt{4'h0}) and \texttt{GEMM} (opcode \texttt{4'h1}).}
\label{fig:type-a}
\end{figure}

\subsection{Type B --- Memory Control Format (MEMCPY)}

\begin{figure}[htb]
\centering
\begin{bytefield}[bitwidth=0.68em, bitheight=2.4em,
                  boxformatting={\centering\sffamily\scriptsize\bfseries}]{64}
  \bitheader[endianness=big,lsb=0]{0,1,7,24,41,58,59,60,63} \\
  \bitbox{4}{opcode\\\scriptsize[63:60]} &
  \bitbox{1}{fd\\\scriptsize[59]} &
  \bitbox{1}{td\\\scriptsize[58]} &
  \bitbox{17}{dest\_addr\\\scriptsize[57:41]} &
  \bitbox{17}{src\_addr\\\scriptsize[40:24]} &
  \bitbox{17}{aux\_addr\\\scriptsize[23:7]} &
  \bitbox{6}{shape\_ptr\\\scriptsize[6:1]} &
  \bitbox{1}{as\\\scriptsize[0]}
\end{bytefield}
\caption{\textbf{Type B} body layout, used by \texttt{MEMCPY} (opcode
\texttt{4'h2}). \texttt{fd}=\texttt{from\_device}, \texttt{td}=\texttt{to\_device},
\texttt{as}=\texttt{async}.}
\label{fig:type-b}
\end{figure}

\subsection{Type C --- Memory Set Format (MEMSET)}

\begin{figure}[htb]
\centering
\begin{bytefield}[bitwidth=0.68em, bitheight=2.4em,
                  boxformatting={\centering\sffamily\scriptsize\bfseries}]{64}
  \bitheader[endianness=big,lsb=0]{0,4,20,36,52,58,60,63} \\
  \bitbox{4}{opcode\\\scriptsize[63:60]} &
  \bitbox{2}{dc\\\scriptsize[59:58]} &
  \bitbox{6}{dest\_addr\\\scriptsize[57:52]} &
  \bitbox{16}{a\_value\\\scriptsize[51:36]} &
  \bitbox{16}{b\_value\\\scriptsize[35:20]} &
  \bitbox{16}{c\_value\\\scriptsize[19:4]} &
  \bitbox{4}{rsv\\\scriptsize[3:0]}
\end{bytefield}
\caption{\textbf{Type C} body layout, used by \texttt{MEMSET} (opcode
\texttt{4'h3}). \texttt{dc}=\texttt{dest\_cache}.}
\label{fig:type-c}
\end{figure}

\subsection{Type D --- Complex Vector Op Format (CVO)}

\vspace{0.3em}
\begin{center}
\begin{bytefield}[bitwidth=0.68em, bitheight=2.4em,
                  boxformatting={\centering\sffamily\scriptsize\bfseries}]{64}
  \bitheader[endianness=big,lsb=0]{0,1,6,22,39,56,60,63} \\
  \bitbox{4}{opcode\\\scriptsize[63:60]} &
  \bitbox{4}{func\\\scriptsize[59:56]} &
  \bitbox{17}{src\_addr\\\scriptsize[55:39]} &
  \bitbox{17}{dst\_addr\\\scriptsize[38:22]} &
  \bitbox{16}{length\\\scriptsize[21:6]} &
  \bitbox{5}{flags\\\scriptsize[5:1]} &
  \bitbox{1}{as\\\scriptsize[0]}
\end{bytefield}
\end{center}
\captionof{figure}{\textbf{Type D} body layout, used by \texttt{CVO} (opcode
\texttt{4'h4}). \texttt{as}=\texttt{async}.}
\label{fig:type-d}
\vspace{0.5em}

\clearpage

% =================================================================
\chapter{INSTRUCTION SET REFERENCE (A-Z)}

This chapter is the normative per-instruction reference. Each entry
follows the \textbf{Intel SDM} format: an Opcode / Instruction /
Description summary table, a Description paragraph, an Operation
pseudocode block, a Flags Affected subsection, and an Exceptions
subsection where applicable.

% =====================================================
% GEMM / GEMV
% =====================================================
\begin{InstructionRef}{GEMM / GEMV}{Matrix-matrix / matrix-vector multiply}

\RefSection{Opcode \& Instruction summary}

\begin{tabularx}{\linewidth}{|c|c|c|X|}
\hline
\rowcolor{skytint}
\textbf{Opcode} & \textbf{Instruction} & \textbf{Op/En} & \textbf{Description} \\
\hline
\texttt{4'h0}\ \ \textit{body} & \texttt{GEMV dest, src, flags, size\_ptr, shape\_ptr, lane} & A & Matrix-vector multiply. 4 GEMV cores consume weights via Weight Streaming. \\
\hline
\texttt{4'h1}\ \ \textit{body} & \texttt{GEMM dest, src, flags, size\_ptr, shape\_ptr, lane} & A & Matrix-matrix multiply. 32$\times$32 systolic array runs Weight Stationary. \\
\hline
\end{tabularx}

\RefSection{Instruction operand encoding}

\begin{tabularx}{\linewidth}{|c|X|}
\hline
\rowcolor{skytint}
\textbf{Op/En} & \textbf{Operand 1 --- Operand 6} \\
\hline
A & \texttt{dest} (Mod ModReg, W) $\vert$ \texttt{src} (Mod ModReg, R) $\vert$ \texttt{flags} (Imm, R) $\vert$ \texttt{size\_ptr} (Ptr, R) $\vert$ \texttt{shape\_ptr} (Ptr, R) $\vert$ \texttt{lane} (Imm, R) \\
\hline
\end{tabularx}

\RefSection{Bit-level encoding (Type A)}

\begin{center}
\begin{bytefield}[bitwidth=0.68em, bitheight=2.1em, boxformatting={\centering\sffamily\scriptsize}]{64}
  \bitheader[endianness=big,lsb=0]{0,3,8,14,20,26,43,60,63} \\
  \bitbox{4}{opcode} &
  \bitbox{17}{dest\_reg\\[-0.1em]\scriptsize[59:43]} &
  \bitbox{17}{src\_addr\\[-0.1em]\scriptsize[42:26]} &
  \bitbox{6}{flags\\\scriptsize[25:20]} &
  \bitbox{6}{size\_ptr\\\scriptsize[19:14]} &
  \bitbox{6}{shape\_ptr\\\scriptsize[13:8]} &
  \bitbox{5}{lane\\\scriptsize[7:3]} &
  \bitbox{3}{rsv}
\end{bytefield}
\end{center}

\RefSection{Description}

Performs a tiled matrix multiplication. \texttt{src} points at the
activation source in L2; \texttt{dest} receives the result. The
dispatcher reads \texttt{shape\_ptr} and \texttt{size\_ptr} from the
Constant Cache to obtain $(M, N, K)$ and tile sizes, then streams a
weight tile into either (a) the 32$\times$32 systolic array for
\texttt{GEMM}, or (b) the 4 parallel GEMV cores for \texttt{GEMV}.
Accumulation is done by the DSP48E2 integer MAC; the BF16 scale
product is applied at post-process.

Setting \texttt{lane = 0} activates all lanes. Values 1--31 restrict
execution to the specified lane count — useful for partial-tile
fallbacks at the trailing edge of a matrix.

\RefSection{Flags field}

\begin{tabularx}{\linewidth}{|c|l|X|}
\hline
\rowcolor{skytint}
\textbf{Bit} & \textbf{Name} & \textbf{Description} \\
\hline
\texttt{[5]}   & \texttt{findemax} & Update the \texttt{e\_max} register for output normalization (used by softmax). \\
\hline
\texttt{[4]}   & \texttt{accm}     & Accumulate into \texttt{dest} (default is overwrite). \\
\hline
\texttt{[3]}   & \texttt{w\_scale} & Apply the weight scale factor during MAC. \\
\hline
\texttt{[2:0]} & reserved          & Must be zero. \\
\hline
\end{tabularx}

\RefSection{Operation}

\begin{lstlisting}[style=pseudo]
DISPATCH gemm_or_gemv:
    shape  <- CONST_CACHE[shape_ptr_addr]       // (M, N, K)
    size   <- CONST_CACHE[size_ptr_addr]        // tile sizes
    FOR each tile (m, n, k) IN iteration_space(shape, size):
        IF opcode == GEMM:
            ARRAY_PRELOAD_WEIGHTS(tile_weights)     // Weight Stationary
            STREAM_ACTIVATIONS(src_addr + offset)
        ELSE:                                       // GEMV
            PARTITION_WEIGHTS_OVER_4_CORES(tile_weights)
            STREAM_ACTIVATIONS(src_addr + offset)
        ENDIF

        IF flags.accm = 1:
            acc <- L2[dest_reg + offset]
        ELSE:
            acc <- 0
        ENDIF

        acc <- acc + INT_MAC(activations, weights) * (S_fmap * S_weight)
        IF flags.findemax = 1:
            EMAX_REG <- max(EMAX_REG, exp(acc))
        ENDIF
        L2[dest_reg + offset] <- acc
    ENDFOR
END DISPATCH
\end{lstlisting}

\RefSection{Flags affected}

\texttt{e\_max} is updated when \texttt{findemax = 1}. No other
architectural flags are affected.

\RefSection{Exceptions}

\begin{itemize}[leftmargin=*, itemsep=0.2ex]
  \item \textbf{\texttt{\#UD} --- undefined instruction.} Raised if
        \texttt{shape\_ptr} or \texttt{size\_ptr} reference an
        uninitialised Constant Cache entry.
  \item \textbf{\texttt{\#RSV}} --- reserved bits \texttt{[2:0]}
        must be zero; decoder asserts on non-zero.
\end{itemize}

\RefSection{Use case}

\begin{itemize}[leftmargin=*, itemsep=0.2ex]
  \item \texttt{GEMM} --- prefill. Q/K/V projection, FFN up/down
        projection over a prompt tile.
  \item \texttt{GEMV} --- autoregressive decoding. Every projection
        in the single-token-per-step loop.
\end{itemize}
\end{InstructionRef}

% =====================================================
% MEMCPY
% =====================================================
\begin{InstructionRef}{MEMCPY}{Move a block between host and device / L2 and L2}

\RefSection{Opcode \& Instruction summary}

\begin{tabularx}{\linewidth}{|c|c|c|X|}
\hline
\rowcolor{skytint}
\textbf{Opcode} & \textbf{Instruction} & \textbf{Op/En} & \textbf{Description} \\
\hline
\texttt{4'h2}\ \ \textit{body} & \texttt{MEMCPY from\_dev, to\_dev, dest, src, aux, shape\_ptr, async} & B & Bulk data move. \texttt{from\_dev}/\texttt{to\_dev} determines direction and path. \\
\hline
\end{tabularx}

\RefSection{Bit-level encoding (Type B)}

\begin{center}
\begin{bytefield}[bitwidth=0.68em, bitheight=2.1em, boxformatting={\centering\sffamily\scriptsize}]{64}
  \bitheader[endianness=big,lsb=0]{0,1,7,24,41,58,59,60,63} \\
  \bitbox{4}{opcode} &
  \bitbox{1}{fd} & \bitbox{1}{td} &
  \bitbox{17}{dest\_addr\\\scriptsize[57:41]} &
  \bitbox{17}{src\_addr\\\scriptsize[40:24]} &
  \bitbox{17}{aux\_addr\\\scriptsize[23:7]} &
  \bitbox{6}{shape\_ptr\\\scriptsize[6:1]} &
  \bitbox{1}{as}
\end{bytefield}
\end{center}

\RefSection{Description}

Moves a shape-defined block between two address spaces. The
\texttt{from\_device}/\texttt{to\_device} pair encodes one of four
source--destination combinations; the decoder expands the pair into
the 8-bit \texttt{data\_route\_e} enum consumed by the back-end DMA
engines. See Table~\ref{tab:data-route}.

\texttt{aux\_addr} carries the host-side DDR offset when the transfer
crosses ACP; for NPU-to-NPU transfers it is ignored.

When \texttt{async = 0} the dispatcher stalls the front end until the
completion fence returns. When \texttt{async = 1} the instruction
retires immediately and completion is surfaced to the host via
\texttt{STAT\_OUT}.

\RefSection{Supported (\texttt{from\_device}, \texttt{to\_device}) combinations}

\begin{tabularx}{\linewidth}{|c|c|X|}
\hline
\rowcolor{skytint}
\textbf{from\_dev} & \textbf{to\_dev} & \textbf{Path} \\
\hline
1 (Host) & 0 (NPU)  & ACP $\to$ L2 cache. Used for weight / input loads. \\
\hline
0 (NPU)  & 1 (Host) & L2 cache $\to$ ACP. Returns output tokens / KV entries. \\
\hline
0 (NPU)  & 0 (NPU)  & Intra-L2 block move (on-device rearrangement). \\
\hline
1 (Host) & 1 (Host) & \textit{Reserved} --- Control Unit raises \texttt{\#UD}. \\
\hline
\end{tabularx}

\RefSection{Operation}

\begin{lstlisting}[style=pseudo]
DISPATCH memcpy:
    route <- ENCODE_ROUTE(from_device, to_device)   // data_route_e
    shape <- CONST_CACHE[shape_ptr_addr]
    CASE route OF
        from_host_to_L2: DMA_ACP_READ(dest_addr, src_addr, aux_addr, shape)
        from_L2_to_host: DMA_ACP_WRITE(dest_addr, src_addr, aux_addr, shape)
        from_L2_to_L1_GEMM: DMA_L1_FWD(dest_addr, src_addr, GEMM, shape)
        from_L2_to_L1_GEMV: DMA_L1_FWD(dest_addr, src_addr, GEMV, shape)
        OTHERWISE: RAISE #UD
    ENDCASE

    IF async = 1:
        SIGNAL fence_id TO STAT_OUT    // fire-and-forget
        RETURN
    ELSE:
        WAIT fence_id
    ENDIF
END DISPATCH
\end{lstlisting}

\RefSection{Exceptions}

\begin{itemize}[leftmargin=*, itemsep=0.2ex]
  \item \textbf{\texttt{\#UD}} on reserved (Host,\,Host) combination.
  \item \textbf{\texttt{\#RSV}} on shape pointer referencing an
        uninitialised Constant Cache entry.
  \item \textbf{\texttt{\#AXI}} if the HP / ACP read returns an
        error response (SLVERR / DECERR). Surfaced via
        \texttt{STAT\_OUT}.
\end{itemize}
\end{InstructionRef}

% =====================================================
% MEMSET
% =====================================================
\begin{InstructionRef}{MEMSET}{Constant Cache write (shape / size / scale)}

\RefSection{Opcode \& Instruction summary}

\begin{tabularx}{\linewidth}{|c|c|c|X|}
\hline
\rowcolor{skytint}
\textbf{Opcode} & \textbf{Instruction} & \textbf{Op/En} & \textbf{Description} \\
\hline
\texttt{4'h3}\ \ \textit{body} & \texttt{MEMSET dest\_cache, dest\_addr, a, b, c} & C & Write three 16-bit values to a Constant Cache entry. \\
\hline
\end{tabularx}

\RefSection{Bit-level encoding (Type C)}

\begin{center}
\begin{bytefield}[bitwidth=0.68em, bitheight=2.1em, boxformatting={\centering\sffamily\scriptsize}]{64}
  \bitheader[endianness=big,lsb=0]{0,4,20,36,52,58,60,63} \\
  \bitbox{4}{opcode} &
  \bitbox{2}{dc\\\scriptsize[59:58]} &
  \bitbox{6}{dest\_addr\\\scriptsize[57:52]} &
  \bitbox{16}{a\_value\\\scriptsize[51:36]} &
  \bitbox{16}{b\_value\\\scriptsize[35:20]} &
  \bitbox{16}{c\_value\\\scriptsize[19:4]} &
  \bitbox{4}{rsv}
\end{bytefield}
\end{center}

\RefSection{Description}

Writes directly into the Constant Cache. This is the \textit{only}
opcode that does not touch the L2 cache. \texttt{dest\_cache} selects
which bank is targeted: 0 = \texttt{fmap\_shape}, 1 =
\texttt{weight\_shape} (2--3 reserved).

The three value slots \texttt{a / b / c} let a single instruction
write an entire $(M, N, K)$ tuple in one shot; this is the idiomatic
use at layer boundaries.

\RefSection{Operation}

\begin{lstlisting}[style=pseudo]
DISPATCH memset:
    CASE dest_cache OF
        0: CONST_CACHE.fmap_shape[dest_addr]   <- (a_value, b_value, c_value)
        1: CONST_CACHE.weight_shape[dest_addr] <- (a_value, b_value, c_value)
        OTHERWISE: RAISE #UD
    ENDCASE
END DISPATCH
\end{lstlisting}

\RefSection{Flags affected} None.

\RefSection{Exceptions}

\begin{itemize}[leftmargin=*, itemsep=0.2ex]
  \item \textbf{\texttt{\#UD}} on reserved \texttt{dest\_cache} values (2--3).
  \item \textbf{\texttt{\#RSV}} on non-zero reserved bits \texttt{[3:0]}.
\end{itemize}
\end{InstructionRef}

% =====================================================
% CVO
% =====================================================
\begin{InstructionRef}{CVO}{Complex Vector Op --- transcendentals \& reductions via SFU}

\RefSection{Opcode \& Instruction summary}

\begin{tabularx}{\linewidth}{|c|c|c|X|}
\hline
\rowcolor{skytint}
\textbf{Opcode} & \textbf{Instruction} & \textbf{Op/En} & \textbf{Description} \\
\hline
\texttt{4'h4}\ \ \textit{body} & \texttt{CVO func, src, dst, length, flags, async} & D & Run an SFU function (exp, sqrt, gelu, sin, cos, reduce, scale, recip) over a length-N vector. \\
\hline
\end{tabularx}

\RefSection{Bit-level encoding (Type D)}

\begin{center}
\begin{bytefield}[bitwidth=0.68em, bitheight=2.1em, boxformatting={\centering\sffamily\scriptsize}]{64}
  \bitheader[endianness=big,lsb=0]{0,1,6,22,39,56,60,63} \\
  \bitbox{4}{opcode} &
  \bitbox{4}{func\\\scriptsize[59:56]} &
  \bitbox{17}{src\_addr\\\scriptsize[55:39]} &
  \bitbox{17}{dst\_addr\\\scriptsize[38:22]} &
  \bitbox{16}{length\\\scriptsize[21:6]} &
  \bitbox{5}{flags\\\scriptsize[5:1]} &
  \bitbox{1}{as}
\end{bytefield}
\end{center}

\RefSection{Description}

Dispatches an SFU (``Scalar Function Unit'') operation on a
\texttt{length}-element BF16 vector. Functions are implemented via a
combination of CORDIC iterations and LUT interpolation; in steady
state the SFU retires one element per cycle.

The \texttt{flags.sub\_emax} bit enables the numerical-stability
pre-subtraction used by softmax ($e^{x - e_{\max}}$). The
\texttt{flags.recip\_scale} bit turns \texttt{CVO\_SCALE} into a
division by flipping the scalar at the input (softmax denominator
divide is implemented this way).

\RefSection{Function codes}

\begin{tabularx}{\linewidth}{|l|c|X|l|}
\hline
\rowcolor{skytint}
\textbf{Function} & \textbf{Code} & \textbf{Semantics} & \textbf{Use} \\
\hline
\texttt{CVO\_EXP}         & \texttt{4'h0} & $e^{x}$                  & Softmax \\
\hline
\texttt{CVO\_SQRT}        & \texttt{4'h1} & $\sqrt{x}$               & RMSNorm \\
\hline
\texttt{CVO\_GELU}        & \texttt{4'h2} & $\mathrm{gelu}(x)$       & FFN nonlinearity \\
\hline
\texttt{CVO\_SIN}         & \texttt{4'h3} & $\sin(x)$                & RoPE \\
\hline
\texttt{CVO\_COS}         & \texttt{4'h4} & $\cos(x)$                & RoPE \\
\hline
\texttt{CVO\_REDUCE\_SUM} & \texttt{4'h5} & $\sum_i x_i$             & Softmax denominator \\
\hline
\texttt{CVO\_SCALE}       & \texttt{4'h6} & $\alpha \cdot x$         & Bias / scale fuse \\
\hline
\texttt{CVO\_RECIP}       & \texttt{4'h7} & $1/x$                    & Softmax / RMSNorm \\
\hline
\texttt{4'h8}--\texttt{4'hF} & --- & \textit{Reserved} & --- \\
\hline
\end{tabularx}

\RefSection{Flags field}

\begin{tabularx}{\linewidth}{|c|l|X|}
\hline
\rowcolor{skytint}
\textbf{Bit} & \textbf{Name} & \textbf{Description} \\
\hline
\texttt{[4]}   & \texttt{sub\_emax}    & Subtract \texttt{e\_max} before computing. \\
\hline
\texttt{[3]}   & \texttt{recip\_scale} & Use the reciprocal of the scalar. \\
\hline
\texttt{[2]}   & \texttt{accm}         & Accumulate into \texttt{dst}. \\
\hline
\texttt{[1:0]} & reserved              & Must be zero. \\
\hline
\end{tabularx}

\RefSection{Operation}

\begin{lstlisting}[style=pseudo]
DISPATCH cvo:
    FOR i FROM 0 TO length - 1 DO
        x <- L2[src_addr + i]
        IF flags.sub_emax = 1:
            x <- x - EMAX_REG
        ENDIF
        CASE func OF
            CVO_EXP:        y <- exp(x)
            CVO_SQRT:       y <- sqrt(x)
            CVO_GELU:       y <- 0.5 * x * (1 + tanh(sqrt(2/pi) * (x + 0.044715*x^3)))
            CVO_SIN:        y <- sin(x)
            CVO_COS:        y <- cos(x)
            CVO_REDUCE_SUM: acc <- acc + x; CONTINUE
            CVO_SCALE:      y <- (flags.recip_scale ? 1/scalar : scalar) * x
            CVO_RECIP:      y <- 1 / x
            OTHERWISE:      RAISE #UD
        ENDCASE

        IF flags.accm = 1:
            y <- y + L2[dst_addr + i]
        ENDIF
        L2[dst_addr + i] <- y
    ENDFOR

    IF func = CVO_REDUCE_SUM:
        L2[dst_addr] <- acc
    ENDIF
    IF async = 1: SIGNAL fence_id; RETURN
    ELSE:         WAIT fence_id
    ENDIF
END DISPATCH
\end{lstlisting}

\RefSection{Fast path}

If the preceding instruction is a \texttt{GEMV} whose \texttt{dest\_reg}
matches a special \textit{bypass tag}, the SFU consumes the output
directly from the GEMV FIFO and skips the L2 round trip.
Figure~\ref{fig:cvo-pipeline} illustrates the data path.

\RefSection{Exceptions}

\begin{itemize}[leftmargin=*, itemsep=0.2ex]
  \item \textbf{\texttt{\#UD}} on reserved function codes (\texttt{4'h8}--\texttt{4'hF}).
  \item \textbf{\texttt{\#RSV}} on non-zero reserved flag bits.
\end{itemize}
\end{InstructionRef}

\section{CVO Pipeline Internal Structure}

\begin{figure}[h]
\centering
\begin{tikzpicture}[
  font=\sffamily\small,
  node distance = 0.3cm and 0.55cm,
  st/.style={rectangle, draw=black, thin, fill=oceanblue!8,
             minimum width=1.7cm, minimum height=0.9cm,
             font=\sffamily\footnotesize, align=center},
  sfu/.style={rectangle, draw=black, thin, fill=accentpurple!8,
              minimum width=1.9cm, minimum height=0.9cm,
              font=\sffamily\footnotesize, align=center},
]
  \node[st] (s1) {Issue};
  \node[st, right=of s1] (s2) {$E_{\max}$\\subtract};
  \node[sfu, right=of s2] (s3) {CORDIC\\iter};
  \node[sfu, right=of s3] (s4) {LUT\\interp.};
  \node[st, right=of s4] (s5) {Scale /\\accm};
  \node[st, right=of s5] (s6) {L2 WB\\or FIFO};

  \draw[arr] (s1) -- (s2);
  \draw[arr] (s2) -- (s3);
  \draw[arr] (s3) -- (s4);
  \draw[arr] (s4) -- (s5);
  \draw[arr] (s5) -- (s6);

  % Bypass arrow from a GEMV FIFO into s1 (fast path, shown dashed)
  \node[blocknode, above=0.5cm of s1, minimum width=2.2cm, fill=accentpurple!8]
    (bp) {GEMV FIFO\\(bypass)};
  \draw[arr, dashed] (bp) -- (s1);
\end{tikzpicture}
\caption{CVO pipeline internals. Dashed red arrow is the fast-path
from a preceding \texttt{GEMV}.}
\label{fig:cvo-pipeline}
\end{figure}

\newpage

% =================================================================
\chapter{MEMORY ROUTING \& ADDRESS SPACES}

\section{Data-Route Encoding (\texttt{data\_route\_e})}

The \texttt{from\_device}/\texttt{to\_device} pair on \texttt{MEMCPY}
expands inside the Control Unit into an 8-bit route enum. Every DMA
controller in the back end consumes this single byte rather than
re-interpret the two encoding bits.

\begin{longtable}{|l|c|c|p{6.5cm}|}
\hline
\rowcolor{skytint}
\textbf{Route symbol} & \textbf{Hex} & \textbf{Direction} & \textbf{Path} \\
\hline
\endfirsthead
\texttt{from\_host\_to\_L2}      & \texttt{8'h01} & Ingress   & Host DDR4 $\to$ L2 cache \\
\hline
\texttt{from\_L2\_to\_host}      & \texttt{8'h10} & Egress    & L2 cache $\to$ Host DDR4 \\
\hline
\texttt{from\_L2\_to\_L1\_GEMM}  & \texttt{8'h12} & Compute   & L2 $\to$ GEMM L1 \\
\hline
\texttt{from\_L2\_to\_L1\_GEMV}  & \texttt{8'h13} & Compute   & L2 $\to$ GEMV L1 \\
\hline
\texttt{from\_L2\_to\_CVO}       & \texttt{8'h14} & Compute   & L2 $\to$ SFU \\
\hline
\texttt{from\_GEMM\_res\_to\_L2} & \texttt{8'h21} & Writeback & GEMM result $\to$ L2 \\
\hline
\texttt{from\_GEMV\_res\_to\_L2} & \texttt{8'h31} & Writeback & GEMV result $\to$ L2 \\
\hline
\texttt{from\_CVO\_res\_to\_L2}  & \texttt{8'h41} & Writeback & SFU result $\to$ L2 \\
\hline
\caption{\texttt{data\_route\_e} encoding and associated datapath.}
\label{tab:data-route}
\end{longtable}

\begin{figure}[h]
\centering
\begin{tikzpicture}[
  font=\sffamily\footnotesize,
  node distance=0.45cm and 0.85cm,
]
  \node[blocknode, minimum width=2.6cm, fill=accentorange!10] (host) {Host DDR4};
  \node[blocknode, right=2.0cm of host, minimum width=2.6cm, fill=accentpurple!8] (l2) {L2 Cache};
  \node[blocknode, above=0.8cm of l2, minimum width=2.6cm, fill=accentpurple!8] (gemm) {GEMM L1};
  \node[blocknode, below=0.8cm of l2, minimum width=2.6cm, fill=accentpurple!8] (gemv) {GEMV L1};
  \node[blocknode, right=2.0cm of l2, minimum width=2.6cm, fill=accentpurple!8]  (sfu)  {SFU};

  \draw[arr]        (host) -- node[above, pos=0.5, fill=white, inner sep=1pt]{\texttt{8'h01}} (l2);
  \draw[arr]        (l2) to [bend right=25] node[below, pos=0.5, fill=white, inner sep=1pt]{\texttt{8'h10}} (host);
  \draw[arr]        (l2)   -- node[right, pos=0.5, fill=white, inner sep=1pt]{\texttt{8'h12}} (gemm);
  \draw[arr]        (l2)   -- node[right, pos=0.5, fill=white, inner sep=1pt]{\texttt{8'h13}} (gemv);
  \draw[arr]        (l2)   -- node[above, pos=0.5, fill=white, inner sep=1pt]{\texttt{8'h14}} (sfu);
  \draw[arr, dashed] (gemm) -- node[left,  pos=0.5, fill=white, inner sep=1pt]{\texttt{8'h21}} (l2);
  \draw[arr, dashed] (gemv) -- node[left,  pos=0.5, fill=white, inner sep=1pt]{\texttt{8'h31}} (l2);
  \draw[arr, dashed] (sfu) to [bend right=25] node[below, pos=0.5, fill=white, inner sep=1pt]{\texttt{8'h41}} (l2);
\end{tikzpicture}
\caption{Visual summary of Table~\ref{tab:data-route}. Solid arrows
are ingress / compute reads; dashed arrows are writeback paths.}
\label{fig:data-route}
\end{figure}

\section{Constant Cache \& Pointer Registers}

The Constant Cache is 64 entries deep and holds shape metadata, tile
sizes, and scale factors. ISA-visible pointers into this cache are
6 bits wide (since $\log_2 64 = 6$).

\begin{longtable}{|l|c|p{8cm}|}
\hline
\rowcolor{skytint}
\textbf{Pointer} & \textbf{Width} & \textbf{Content} \\
\hline
\endfirsthead
\texttt{shape\_ptr\_addr}  & 6 & $(M, N, K)$ tuple for the tile dimension. \\
\hline
\texttt{size\_ptr\_addr}   & 6 & Tile sizes, loop bounds, stride. \\
\hline
\texttt{ptr\_addr\_t}      & 6 & Generic index into the 64-entry Constant Cache. \\
\hline
\end{longtable}

Pointer entries are populated by \texttt{MEMSET}. A canonical layer
prologue is:

\begin{lstlisting}[style=pseudo]
MEMSET fmap_shape,   idx=0, a=M, b=N, c=K
MEMSET weight_shape, idx=0, a=scale, b=zero_point, c=0
GEMM   dest=..., src=..., flags={w_scale=1}, size_ptr=0, shape_ptr=0, lane=0
\end{lstlisting}

\section{Address Space}

\begin{longtable}{|l|c|p{9cm}|}
\hline
\rowcolor{skytint}
\textbf{Field} & \textbf{Width} & \textbf{Address space} \\
\hline
\endfirsthead
\texttt{dest\_addr} / \texttt{src\_addr} & 17 & $2^{17}$ entries = 128\,K, indexed by L2 cache block. \\
\hline
\texttt{aux\_addr}                       & 17 & MEMCPY auxiliary address (e.g., host DDR offset in units of the block size). \\
\hline
\end{longtable}

With an on-device block size of 128 bits (16 B) on KV260, the 17-bit
field yields a \textbf{2\,MB} linear L2 address space.

\begin{CautionBox}
\texttt{aux\_addr} is measured in \textbf{block units}, not bytes.
A host-side byte offset must be divided by 16 before loading into
\texttt{aux\_addr}.
\end{CautionBox}

\newpage

% =================================================================
\chapter{DATA FLOW \& COMPLETION}

\section{Per-Instruction Dataflow Summary}

\begin{figure}[h]
\centering
\resizebox{\textwidth}{!}{%
\begin{tikzpicture}[
  font=\sffamily\footnotesize,
  node distance=0.3cm and 0.6cm,
  dfstep/.style={rectangle, draw=black, thin, fill=oceanblue!8,
               minimum width=2.6cm, minimum height=0.8cm,
               font=\sffamily\footnotesize, align=center},
]
  % GEMM path
  \node[dfstep] (g1) {Const\,Cache\\lookup};
  \node[dfstep, right=of g1] (g2) {Weight\\preload};
  \node[dfstep, right=of g2] (g3) {Activation\\stream};
  \node[dfstep, right=of g3] (g4) {Systolic\\accumulate};
  \node[dfstep, right=of g4] (g5) {Post-proc\\(scale,bias)};
  \node[dfstep, right=of g5] (g6) {WB L2};
  \node[above=0.15cm of g1, font=\sffamily\bfseries\color{oceanblue}] {GEMM};
  \foreach \a/\b in {g1/g2, g2/g3, g3/g4, g4/g5, g5/g6} \draw[arr] (\a)--(\b);

  % GEMV path
  \node[dfstep, below=0.55cm of g1] (v1) {Const\,Cache\\lookup};
  \node[dfstep, right=of v1] (v2) {Row-part.\\weight};
  \node[dfstep, right=of v2] (v3) {4 cores\\32-MAC};
  \node[dfstep, right=of v3] (v4) {5-stage\\reduce};
  \node[dfstep, right=of v4] (v5) {Post-proc};
  \node[dfstep, right=of v5] (v6) {WB L2\\or SFU};
  \node[above=0.15cm of v1, font=\sffamily\bfseries\color{oceanblue}] {GEMV};
  \foreach \a/\b in {v1/v2, v2/v3, v3/v4, v4/v5, v5/v6} \draw[arr] (\a)--(\b);

  % MEMCPY path
  \node[dfstep, below=0.55cm of v1] (m1) {Resolve\\route};
  \node[dfstep, right=of m1] (m2) {Shape\\lookup};
  \node[dfstep, right=of m2] (m3) {DMA\\issue};
  \node[dfstep, right=of m3] (m4) {Bus\\transfer};
  \node[dfstep, right=of m4] (m5) {Fence\\register};
  \node[dfstep, right=of m5] (m6) {STAT\_OUT\\report};
  \node[above=0.15cm of m1, font=\sffamily\bfseries\color{oceanblue}] {MEMCPY};
  \foreach \a/\b in {m1/m2, m2/m3, m3/m4, m4/m5, m5/m6} \draw[arr] (\a)--(\b);

  % CVO path
  \node[dfstep, below=0.55cm of m1] (c1) {Func\\decode};
  \node[dfstep, right=of c1] (c2) {L2 read\\src};
  \node[dfstep, right=of c2] (c3) {$E_{\max}$\\sub};
  \node[dfstep, right=of c3] (c4) {CORDIC/\\LUT};
  \node[dfstep, right=of c4] (c5) {Scale/\\accm};
  \node[dfstep, right=of c5] (c6) {WB L2\\or FIFO};
  \node[above=0.15cm of c1, font=\sffamily\bfseries\color{oceanblue}] {CVO};
  \foreach \a/\b in {c1/c2, c2/c3, c3/c4, c4/c5, c5/c6} \draw[arr] (\a)--(\b);
\end{tikzpicture}%
}
\caption{Per-instruction dataflow summary. Each row is one
instruction's journey through the NPU from dispatch to writeback.}
\label{fig:dataflow-summary}
\end{figure}

\section{Hazards and Completion Tracking}

\begin{figure}[h]
\centering
\resizebox{\textwidth}{!}{%
\begin{tikzpicture}[
  font=\sffamily\footnotesize,
  node distance=0.3cm and 0.6cm,
]
  \node[blocknode, minimum width=3cm] (in) {Incoming \uop};
  \node[blocknode, right=1.3cm of in, minimum width=3.4cm, fill=accentorange!10]
    (ah) {Address hazard check\\\scriptsize RAW on dest/src};
  \node[blocknode, right=1.3cm of ah, minimum width=3.4cm, fill=accentorange!10]
    (rh) {Resource hazard check\\\scriptsize GEMM/GEMV/SFU};
  \node[blocknode, right=1.3cm of rh, minimum width=2.4cm, fill=accentred!8] (ex) {Execute};
  \node[scheduler, below=0.55cm of rh, minimum width=7cm]
    (fence) {Fence tracker --- \texttt{fsmout\_npu\_stat\_collector}};
  \node[blocknode, right=0.6cm of ex, minimum width=2.5cm, fill=accentpurple!8] (stat) {STAT\_OUT};

  \draw[arr] (in) -- (ah);
  \draw[arr] (ah) -- (rh);
  \draw[arr] (rh) -- (ex);
  \draw[arr] (ex.south) -- ++(0,-0.4cm) -| (fence.east);
  \draw[arr] (fence.east) -| (stat.south);
\end{tikzpicture}%
}
\caption{Hazard pipeline. Each \uop passes an address-hazard and a
resource-hazard check before executing; \texttt{async} completion is
reported via the Fence tracker into \texttt{STAT\_OUT}.}
\label{fig:hazard}
\end{figure}

\section{Asynchronous Completion Protocol}

For \texttt{async = 1} instructions (MEMCPY or CVO), the dispatcher
retires immediately and the host tracks completion through
\texttt{STAT\_OUT}. Figure~\ref{fig:async-seq} shows the wire-level
sequence.

\begin{figure}[h]
\centering
\begin{tikzpicture}[
  font=\sffamily\small,
  node distance=0.2cm and 1.2cm,
  agent/.style={rectangle, draw=black, thin, fill=oceanblue!8,
                minimum width=2.4cm, minimum height=0.7cm, font=\sffamily\small, align=center},
]
  \node[agent] (host) {Host driver};
  \node[agent, right=3.5cm of host] (ctrl) {Control Unit};
  \node[agent, right=3.5cm of ctrl] (be)   {Back end (DMA/SFU)};

  \coordinate (h0) at ($(host.south) + (0,-0.2cm)$);
  \coordinate (c0) at ($(ctrl.south) + (0,-0.2cm)$);
  \coordinate (b0) at ($(be.south)   + (0,-0.2cm)$);

  \draw[black!50, thin] (h0) -- ++(0,-5.5cm);
  \draw[black!50, thin] (c0) -- ++(0,-5.5cm);
  \draw[black!50, thin] (b0) -- ++(0,-5.5cm);

  % Message arrows
  \draw[arrfat] ($(h0)+(0,-0.5cm)$) -- node[above, font=\sffamily\scriptsize]{\texttt{CMD\_IN} 64\,b (async=1)} ($(c0)+(0,-0.5cm)$);
  \draw[arrfat] ($(c0)+(0,-1.2cm)$) -- node[above, font=\sffamily\scriptsize]{\uop + \texttt{fence\_id}} ($(b0)+(0,-1.2cm)$);
  \draw[arrfat] ($(c0)+(0,-1.9cm)$) -- node[above, font=\sffamily\scriptsize]{ack (retire to host)} ($(h0)+(0,-1.9cm)$);
  \draw[arrdim, dashed] ($(b0)+(0,-2.6cm)$) -- node[below, font=\sffamily\scriptsize]{(bus transfer \ldots\ latency)} ($(c0)+(0,-2.6cm)$);
  \draw[arrfat] ($(b0)+(0,-3.8cm)$) -- node[above, font=\sffamily\scriptsize]{fence done} ($(c0)+(0,-3.8cm)$);
  \draw[arrfat] ($(c0)+(0,-4.5cm)$) -- node[above, font=\sffamily\scriptsize]{\texttt{STAT\_OUT} update} ($(h0)+(0,-4.5cm)$);
\end{tikzpicture}
\caption{Asynchronous completion sequence. The host can issue the
next instruction as soon as it receives the retire ack; true
completion arrives later via \texttt{STAT\_OUT}.}
\label{fig:async-seq}
\end{figure}

\section{Fence Tracker State Machine}

The fence tracker is a small FSM inside the Dispatcher with one
instance per outstanding \texttt{async} instruction.  Each instance
owns a single \texttt{fence\_id} slot in \texttt{STAT\_OUT}.

\begin{figure}[h]
\centering
\begin{tikzpicture}[
  font=\sffamily\small,
  node distance = 1.8cm and 2.8cm,
  state/.style = {circle, draw=black, thin, minimum size=1.7cm, align=center,
                  font=\sffamily\footnotesize, inner sep=1pt, fill=oceanblue!8},
  stateh/.style= {circle, draw=black, thin, minimum size=1.7cm, align=center,
                  font=\sffamily\footnotesize, inner sep=1pt, fill=accentred!8},
  edgelbl/.style={font=\sffamily\scriptsize, fill=white, inner sep=1pt},
]
  \node[state]  (idle) {\textbf{IDLE}\\\scriptsize slot free};
  \node[stateh, right=of idle] (track) {\textbf{TRACKING}\\\scriptsize \uop in\\\scriptsize flight};
  \node[state, right=of track] (done) {\textbf{DONE}\\\scriptsize fence\\\scriptsize retired};

  % Transitions
  \draw[arr, bend left=15]   (idle.north east)
        to node[edgelbl, above] {\texttt{async} \uop dispatched}
        (track.north west);
  \draw[arr, bend left=15]   (track.north east)
        to node[edgelbl, above] {back-end \texttt{fence\_done}}
        (done.north west);
  \draw[arr, bend left=30]   (done.south west)
        to node[edgelbl, below] {host reads \texttt{STAT\_OUT}}
        (idle.south east);

  % Self-loop on TRACKING
  \draw[arr] (track.north)
        .. controls +(150:1.0) and +(30:1.0) ..
        node[edgelbl, above] {partial progress (no state change)}
        (track.north);

  % Reset from any state on reset (abstracted)
  \node[font=\sffamily\scriptsize\color{graytext}, below=2.2cm of track, align=center]
    {Global reset returns every tracker slot to \textbf{IDLE}.\\
     There are 16 slots in hardware --- one per outstanding \texttt{fence\_id}.};
\end{tikzpicture}
\caption{Per-slot fence-tracker state machine.  A slot transitions
IDLE$\to$TRACKING the moment an \texttt{async} \uop is accepted by the
back end, and TRACKING$\to$DONE when the back end asserts
\texttt{fence\_done}.  The host reading \texttt{STAT\_OUT} returns the
slot to IDLE.}
\label{fig:fence-fsm}
\end{figure}

\newpage

% =================================================================
\chapter{PROGRAMMING GUIDE}

This chapter walks through representative instruction streams that a
host driver emits.

\section{Example 1 --- Single GEMM tile (Q projection for one prompt tile)}

\textbf{Goal:} compute $Q = X W_Q$ for a tile of shape
$(M, N, K) = (64, 256, 256)$.

\begin{lstlisting}[style=sv]
// 1. Preload shape + size pointers
MEMSET fmap_shape,   idx=0, a=64,  b=256, c=256     // (M, N, K)
MEMSET weight_shape, idx=0, a=scale_wq, b=zp_wq, c=0

// 2. Load weights if not already streaming
MEMCPY from_device=1, to_device=0, dest=L2_W_BASE,
       src=HOST_WQ_OFFSET, aux=HOST_BASE,
       shape_ptr=0, async=0

// 3. Dispatch the GEMM
GEMM   dest_reg=L2_Q_BASE, src_addr=L2_X_BASE,
       flags={findemax=0, accm=0, w_scale=1},
       size_ptr=0, shape_ptr=0, lane=0                // lane=0 -> all lanes
\end{lstlisting}

\section{Example 2 --- Autoregressive step (GEMV + softmax + RoPE)}

\begin{lstlisting}[style=sv]
// --- Q, K, V projection ---
GEMV dest=L2_Q, src=L2_X, flags=w_scale, size_ptr=0, shape_ptr=0
GEMV dest=L2_K, src=L2_X, flags=w_scale, size_ptr=1, shape_ptr=1
GEMV dest=L2_V, src=L2_X, flags=w_scale, size_ptr=2, shape_ptr=2

// --- RoPE rotation on Q and K ---
CVO  func=CVO_COS, src=L2_Q, dst=L2_Q_R, length=D_HEAD, flags=0
CVO  func=CVO_SIN, src=L2_K, dst=L2_K_R, length=D_HEAD, flags=0

// --- Attention score + softmax ---
GEMV dest=L2_SCORES, src=L2_Q_R, flags={findemax=1, w_scale=0},
     size_ptr=3, shape_ptr=3                        // Q . K^T
CVO  func=CVO_EXP, src=L2_SCORES, dst=L2_SCORES,
     length=SEQ_LEN, flags=sub_emax
CVO  func=CVO_REDUCE_SUM, src=L2_SCORES, dst=L2_DENOM,
     length=SEQ_LEN, flags=0
CVO  func=CVO_SCALE, src=L2_SCORES, dst=L2_SCORES,
     length=SEQ_LEN, flags=recip_scale              // divide by denom

// --- Attention-weighted V gather ---
GEMV dest=L2_CTX, src=L2_SCORES, flags=w_scale,
     size_ptr=4, shape_ptr=4                        // V . P
\end{lstlisting}

\section{Example 3 --- Weight hot-swap (async prefetch overlapped with compute)}

\begin{lstlisting}[style=sv]
// Fire and forget -- continues to the next instruction immediately
MEMCPY from_device=1, to_device=0,
       dest=L2_W_NEXT_BASE,
       src=HOST_W_LAYER_NP1_OFFSET,
       aux=HOST_BASE,
       shape_ptr=5, async=1

// ... overlap with layer N GEMMs ...
GEMM   dest=..., src=..., ...
GEMM   dest=..., src=..., ...

// Before starting layer N+1, the scheduler knows the MEMCPY fence has
// retired, so no explicit barrier is required.
\end{lstlisting}

\section{Performance Rules of Thumb}

\begin{itemize}[label={\color{oceanblue}\textbullet}, leftmargin=*, itemsep=0.3ex]
    \item Amortize \texttt{MEMSET} --- one at the top of each layer is
          enough; avoid per-tile MEMSETs.
    \item Prefer \texttt{async = 1} for MEMCPY when the transfer is
          followed by $\ge 2$ compute instructions that do not depend
          on its result. The scheduler's fence tracker will stall
          automatically when the consumer arrives.
    \item Keep \texttt{parallel\_lane = 0} unless you are deliberately
          executing a trailing partial tile — the dispatcher's
          resource allocator is more efficient when given the full
          lane set.
    \item Fuse softmax: on the same SFU dispatch, chain
          \texttt{CVO\_EXP} $\to$ \texttt{CVO\_REDUCE\_SUM} $\to$
          \texttt{CVO\_SCALE}\allowbreak(\texttt{recip\_scale}) so the
          pipeline drains once rather than three times.
    \item If a GEMV feeds directly into a CVO, pre-tag the GEMV
          \texttt{dest} with the SFU-bypass marker so the L2 round
          trip is skipped.
\end{itemize}

\section{Roofline Model}

Figure~\ref{fig:roofline} plots the classical roofline for pccx v002:
the ceiling is the $800$\,GMAC/s peak compute of 2000 DSP48E2 slices
running at $400$\,MHz in $9 \times 4$ packed mode, and the slope is
the $12.8$\,GB/s DDR4 bandwidth delivered through the 4 AXI-HP ports.
The three operators the ISA exposes land in distinct regimes:

\begin{itemize}[label={\color{oceanblue}\textbullet}, leftmargin=*, itemsep=0.3ex]
    \item \textbf{GEMM} (prefill) — weight reuse of $\sim 256 \times$
          inside one tile pushes arithmetic intensity above the ridge
          point; the systolic array runs compute-bound.
    \item \textbf{GEMV} (decode) — every weight is used once per
          token, so AI is $\sim 0.5$ MAC/byte and throughput tracks
          the memory slope.  Speculative decoding and tied-embedding
          reuse are the only ways to cross the ridge.
    \item \textbf{CVO} — dominated by the activation/scale memory
          traffic; runs left of the ridge but is a small fraction of
          total cycles.
\end{itemize}

\begin{figure}[h]
\centering
\begin{tikzpicture}
  \begin{loglogaxis}[
    width=10.5cm,
    height=7.2cm,
    xlabel={Arithmetic intensity (MAC / byte)},
    ylabel={Achieved throughput (GMAC/s)},
    xmin=0.1, xmax=1000,
    ymin=1,   ymax=2000,
    grid=both,
    major grid style={draw=graytext!30},
    minor grid style={draw=graytext!15},
    axis line style={draw=graytext},
    tick label style={font=\sffamily\footnotesize, color=graytext},
    label style={font=\sffamily\small},
    legend style={font=\sffamily\footnotesize, at={(0.98,0.05)}, anchor=south east,
                  draw=graytext!40, fill=white},
    legend cell align=left,
  ]
    % Memory slope: y = 12.8 GB/s * x (MAC/byte) -> treat MAC as 1 op each
    % and bandwidth as 12.8 byte/sec (in GB/s).  So y_max_mem = BW * AI.
    \addplot[domain=0.1:100, samples=2, semithick, color=black!70]
      { 12.8 * x };
    \addlegendentry{Memory bound (12.8\,GB/s)}

    % Compute ceiling at 800 GMAC/s
    \addplot[domain=0.1:1000, samples=2, semithick, color=black, dashed]
      { 800 };
    \addlegendentry{Compute peak (800\,GMAC/s)}

    % Operating points
    \addplot[only marks, mark=*, mark size=2.7pt, color=black]
      coordinates { (0.5, 6.4) };
    \addlegendentry{GEMV (AI $\approx$ 0.5)}

    \addplot[only marks, mark=square*, mark size=2.7pt, color=black!70]
      coordinates { (256, 800) };
    \addlegendentry{GEMM (AI $\approx$ 256)}

    \addplot[only marks, mark=triangle*, mark size=3pt, color=black!50]
      coordinates { (2, 25.6) };
    \addlegendentry{CVO (AI $\approx$ 2)}

    % Ridge point annotation
    \addplot[dashed, semithick, color=black!40] coordinates { (62.5, 1) (62.5, 800) };
    \node[font=\sffamily\scriptsize\color{graytext}, anchor=south west]
      at (axis cs: 62.5, 1.2) {ridge: 62.5 MAC/B};
  \end{loglogaxis}
\end{tikzpicture}
\caption{Roofline model for pccx v002.  GEMM sits at the compute
ceiling; GEMV is firmly memory-bound.  CVO is off-critical-path but
also bandwidth-limited.}
\label{fig:roofline}
\end{figure}

\newpage

% =================================================================
\appendix

\chapter{OPCODE MAP}

\section{Single-Byte Opcode Quick Reference}

\begin{longtable}{|c|l|l|p{5.5cm}|c|}
\hline
\rowcolor{skytint}
\textbf{opcode[3:0]} & \textbf{Mnemonic} & \textbf{Body layout} & \textbf{Primary back-end} & \textbf{Family} \\
\hline
\endfirsthead
\texttt{4'h0} & GEMV   & Type A & 4 $\times$ 32-MAC streaming cores.   & Compute \\
\hline
\texttt{4'h1} & GEMM   & Type A & 32 $\times$ 32 systolic array.       & Compute \\
\hline
\texttt{4'h2} & MEMCPY & Type B & ACP + HP DMA engines.                & Memory \\
\hline
\texttt{4'h3} & MEMSET & Type C & Constant Cache (bypasses L2).        & Memory \\
\hline
\texttt{4'h4} & CVO    & Type D & SFU (BF16 CORDIC + LUT).             & SFU \\
\hline
\texttt{4'h5}--\texttt{4'hF} & --- & --- & \textit{Undefined}. & --- \\
\hline
\end{longtable}

\chapter{DATA-ROUTE ENUM REFERENCE}

See Table~\ref{tab:data-route} on the corresponding page for the
full text; this appendix reproduces only the 8-bit hex values for
quick lookup.

\begin{longtable}{|c|l|}
\hline
\rowcolor{skytint}
\textbf{Hex} & \textbf{Route symbol} \\
\hline
\endfirsthead
\texttt{8'h01} & \texttt{from\_host\_to\_L2} \\
\hline
\texttt{8'h10} & \texttt{from\_L2\_to\_host} \\
\hline
\texttt{8'h12} & \texttt{from\_L2\_to\_L1\_GEMM} \\
\hline
\texttt{8'h13} & \texttt{from\_L2\_to\_L1\_GEMV} \\
\hline
\texttt{8'h14} & \texttt{from\_L2\_to\_CVO} \\
\hline
\texttt{8'h21} & \texttt{from\_GEMM\_res\_to\_L2} \\
\hline
\texttt{8'h31} & \texttt{from\_GEMV\_res\_to\_L2} \\
\hline
\texttt{8'h41} & \texttt{from\_CVO\_res\_to\_L2} \\
\hline
\end{longtable}

\chapter{CVO FUNCTION-CODE REFERENCE}

\begin{longtable}{|l|c|l|l|}
\hline
\rowcolor{skytint}
\textbf{Function} & \textbf{Code} & \textbf{Semantics} & \textbf{Use} \\
\hline
\endfirsthead
\texttt{CVO\_EXP}         & \texttt{4'h0} & $e^{x}$             & Softmax numerator \\
\hline
\texttt{CVO\_SQRT}        & \texttt{4'h1} & $\sqrt{x}$          & RMSNorm denominator \\
\hline
\texttt{CVO\_GELU}        & \texttt{4'h2} & $\mathrm{gelu}(x)$  & FFN nonlinearity \\
\hline
\texttt{CVO\_SIN}         & \texttt{4'h3} & $\sin(x)$           & RoPE \\
\hline
\texttt{CVO\_COS}         & \texttt{4'h4} & $\cos(x)$           & RoPE \\
\hline
\texttt{CVO\_REDUCE\_SUM} & \texttt{4'h5} & $\sum_i x_i$        & Softmax denominator \\
\hline
\texttt{CVO\_SCALE}       & \texttt{4'h6} & $\alpha \cdot x$    & Bias / scale fuse \\
\hline
\texttt{CVO\_RECIP}       & \texttt{4'h7} & $1/x$               & Softmax / RMSNorm \\
\hline
\texttt{4'h8}--\texttt{4'hF} & --- & \textit{Reserved} & --- \\
\hline
\end{longtable}

\chapter{GLOSSARY}

\begin{longtable}{|>{\raggedright\arraybackslash}p{3.8cm}|p{11.8cm}|}
\hline
\rowcolor{skytint}
\textbf{Term} & \textbf{Definition} \\
\hline
\endfirsthead

\rowcolor{skytint}
\textbf{Term} & \textbf{Definition} \\
\hline
\endhead

\texttt{accm} & GEMM/GEMV flag (bit [4] of the flags field) --- when set,
the result is \textbf{added} to the destination; when clear, the
destination is overwritten. \\
\hline
ACP & Accelerator Coherency Port --- a cache-coherent AXI port on the
Zynq UltraScale+ PS side that the NPU uses for host-visible DMA. \\
\hline
\uop & Micro-operation: a decoded, back-end-specific packet emitted
by the Dispatcher.  One architectural instruction fans out to one or
more \uops. \\
\hline
Async / \texttt{async} & MEMCPY / CVO flag (bit [0] of the bit-0 flag)
--- when set, the Control Unit retires the instruction immediately and
reports completion through \texttt{STAT\_OUT} instead of blocking. \\
\hline
AXI-Lite & Lightweight subset of the ARM AMBA AXI4 protocol used here
for the 64-bit \texttt{CMD\_IN} / \texttt{STAT\_OUT} control path.
Single-beat, non-burst. \\
\hline
AXI-HP & High-Performance AXI port on the KV260's PS --- the NPU claims
four 128-bit HP ports for bulk data movement between DDR4 and L2. \\
\hline
BF16 & ``Brain floating-point'' half-precision format: 1 sign / 8
exponent / 7 mantissa bits.  Used for activations, KV cache, and SFU
intermediates.  Chosen for its IEEE-float-compatible exponent range
without the FP16 precision cliff. \\
\hline
BFP & Block Floating Point --- an activation-level numeric format
where a group of INT8 mantissas shares one BF16 exponent
($S_{\mathrm{fmap}}$).  Folds the MAC product into a single post-
process scalar multiply. \\
\hline
BRAM / URAM & Xilinx on-chip RAM primitives: BRAM is 36\,Kb / block,
URAM is 288\,Kb / block.  v002 uses URAM as the unified L2 backing
store and BRAM for per-core L1. \\
\hline
CDC & Clock Domain Crossing --- the boundary between the 250\,MHz
AXI domain and the 400\,MHz compute domain.  Crossed by XPM async
FIFOs; see Figure~\ref{fig:cdc}. \\
\hline
CORDIC & ``COordinate Rotation DIgital Computer'' --- an iterative
shift-and-add algorithm that evaluates trigonometric and hyperbolic
functions without a hardware multiplier.  Used inside the SFU for
\texttt{CVO\_SIN}, \texttt{CVO\_COS}, \texttt{CVO\_EXP},
\texttt{CVO\_SQRT}. \\
\hline
CVO & Complex Vector Op --- the family of non-linear scalar/vector
activations (softmax, GELU, RMSNorm, RoPE) dispatched to the SFU.
See \texttt{CVO\_*} function codes in Appendix~C. \\
\hline
\texttt{dest\_cache} (\texttt{dc}) & MEMSET operand: 2-bit selector for
which Constant Cache bank receives the write
(0=\texttt{fmap\_shape}, 1=\texttt{weight\_shape}, 2--3 reserved). \\
\hline
DSP48E2 & Xilinx UltraScale+ DSP slice: $30 \times 18$ signed multiplier
plus 48-bit accumulator.  v002 drives it in $9 \times 4$ packed W4A8
mode for two MACs / cycle per slice. \\
\hline
Dispatcher & Control-unit block that routes decoded \uops to the
back-ends (GEMM, GEMV, SFU, DMA).  Also runs the fence tracker. \\
\hline
$E_{\max}$ & Maximum exponent found during a softmax numerator pass;
subtracted before the \texttt{CVO\_EXP} to keep the sum finite. \\
\hline
Fence tracker & State machine inside the Dispatcher that tracks
in-flight \texttt{async} instructions and writes
\texttt{STAT\_OUT} when the back-end fence retires; see
Figure~\ref{fig:fence-fsm}. \\
\hline
\texttt{findemax} & GEMV flag (bit [5]) --- update the $E_{\max}$
register during this MAC epoch; consumed by a following
\texttt{CVO\_EXP}. \\
\hline
GEMM & General Matrix-Matrix Multiply (prefill-dominant).  v002 uses a
$32 \times 32$ systolic array with a cascade break at row 16. \\
\hline
GEMV & General Matrix-Vector Multiply (decode-dominant).  v002 has
4 GEMV cores each with 32-MAC streaming and a 5-stage reduction tree;
see Figure~\ref{fig:gemv-tree}. \\
\hline
HP port & See AXI-HP. \\
\hline
INT4 / INT8 & 4-bit / 8-bit signed integer.  INT4 is the weight
quantization target (W4); INT8 is the activation quantization target
(A8) in the W4A8 scheme. \\
\hline
KV cache & Key / Value cache --- a per-layer BF16 buffer in L2 that
stores attention keys and values across tokens during autoregressive
decode. \\
\hline
Lane / \texttt{parallel\_lane} & GEMM/GEMV operand (5 bits) that
restricts execution to a subset of the 32 lanes.  \texttt{lane = 0}
activates all lanes. \\
\hline
L1 / L2 & L1 = per-core weight buffer backed by URAM FIFOs.  L2 =
unified 1.75\,MB central cache shared by GEMM, GEMV, SFU, DMA. \\
\hline
LUT (Xilinx) & 6-input lookup table, the basic configurable-logic
cell.  v002 uses LUTs for GEMV dequantize pre-computation and for
reduction tree stages 2--5. \\
\hline
MEMCPY & DMA move between two address spaces.  Can be synchronous
(blocking) or \texttt{async = 1} (fire-and-forget). \\
\hline
MEMSET & Constant Cache scalar write.  Used to push
($M$, $N$, $K$) and scale factors to pointer indices before a
GEMM/GEMV. \\
\hline
opcode & 4-bit field in bits [63:60] of every v002 instruction.  See
Appendix~A for the full map. \\
\hline
PL / PS & Programmable Logic (FPGA fabric) / Processing System (ARM
cores) --- the two halves of the Zynq UltraScale+ SoC on KV260. \\
\hline
\texttt{recip\_scale} & Reciprocal of the softmax denominator;
supplied as an operand to the fused \texttt{CVO\_SCALE} stage to
complete the softmax division in one shot. \\
\hline
RoPE & Rotary Position Embedding --- applies $\sin$ / $\cos$
rotations to Q and K vectors before attention.  Expressed in pccx as
paired \texttt{CVO\_SIN} / \texttt{CVO\_COS} dispatches. \\
\hline
\texttt{shape\_ptr} / \texttt{size\_ptr} & 6-bit operand fields
pointing into the Constant Cache for the ($M$, $N$, $K$) tuple and
for the tile-size tuple.  Populated by MEMSET. \\
\hline
SFU & Special Function Unit --- the single-core back-end that
executes CVO instructions.  Mixes CORDIC iterations with LUT-based
interpolation.  BF16 datapath. \\
\hline
\texttt{src\_addr} / \texttt{dest\_addr} & 17-bit operand fields
addressing 128-bit blocks in L2 (2\,MB linear). \\
\hline
STAT\_OUT & Host-side status register written by the Dispatcher once
an \texttt{async} instruction's fence retires.  AXI-Lite read-only. \\
\hline
Systolic array & Regular grid of PEs used for GEMM.  v002 pins
weights inside PEs (weight-stationary) while activations march
left-to-right and partial sums march top-to-bottom.  See
Figure~\ref{fig:systolic}. \\
\hline
Weight Stationary / Weight Streaming & Two systolic schedules.
\textbf{Stationary} (GEMM) reuses the same weight across many
activations; \textbf{Streaming} (GEMV) streams fresh weights every
cycle because each weight is used once per token. \\
\hline
W4A8 & The pccx v002 quantization scheme: 4-bit weights, 8-bit
activations, BF16 scale factors applied at post-process. \\
\hline
XPM & Xilinx Parameterized Macro --- vendor-provided primitives for
FIFOs, CDC, memory etc.  The NPU uses \texttt{xpm\_fifo\_async} in
independent-clock mode for the PS~$\leftrightarrow$~PL crossing. \\
\hline
\end{longtable}

\newpage

% =================================================================
\chapter{REVISION HISTORY}

\begin{longtable}{|l|l|p{9.5cm}|}
\hline
\rowcolor{skytint}
\textbf{Revision} & \textbf{Date} & \textbf{Notes} \\
\hline
\endfirsthead
v002 & 2026-04-21 &
Initial Software Developer's Manual release.  Intel-SDM-style
per-instruction reference blocks, bit-level Type A/B/C/D
\texttt{bytefield} diagrams, CUDA-grade TikZ architectural
illustrations (system context, NPU top-level, memory hierarchy,
URAM allocation, CDC FIFO, DSP48E2 datapath, systolic array, GEMV
streaming, decode/dispatch, hazard tracker, async completion
sequence, CVO pipeline).  Numbers mirror \texttt{isa\_pkg.sv} in
the RTL repo. \\
\hline
v002-r2 & 2026-04-22 &
Intel-SDM polish pass.  Added wire-level AXI-HP read-handshake
timing diagram (Figure~\ref{fig:axi-timing}), GEMV 5-stage
reduction tree (Figure~\ref{fig:gemv-tree}), fence-tracker state
machine (Figure~\ref{fig:fence-fsm}), and a roofline performance
model (Figure~\ref{fig:roofline}).  New Appendix~D glossary
(~40 terms). \\
\hline
\end{longtable}

\end{document}
